\documentclass[11pt,a4paper]{article}

\usepackage[margin=2.5cm]{geometry}
\usepackage{amsmath,amssymb,amsthm}
\usepackage{booktabs}
\usepackage{longtable}
\usepackage{array}
\usepackage[shortlabels]{enumitem}
\usepackage{hyperref}
\usepackage{xcolor}
\usepackage{listings}
\usepackage{fancyvrb}
\usepackage{parskip}
\usepackage{microtype}
\usepackage{tabularx}
\usepackage{authblk}
\usepackage{abstract}
\sloppy
\emergencystretch=3em
\hbadness=10000

% Theorem environments (academic-paper standard).
\newtheorem{theorem}{Theorem}[section]
\newtheorem{lemma}[theorem]{Lemma}
\newtheorem{proposition}[theorem]{Proposition}
\newtheorem{corollary}[theorem]{Corollary}
\theoremstyle{definition}
\newtheorem{definition}[theorem]{Definition}
\newtheorem{construction}[theorem]{Construction}
\theoremstyle{remark}
\newtheorem{remark}[theorem]{Remark}

% Role tag (kept for cross-reference with implementation notes).
\newcommand{\role}[1]{\textsc{\small[#1]}\,}

% Subsection-level phase header.
\newcommand{\phase}[2]{%
  \subsection[#1]{\texorpdfstring{\textbf{#1} --- #2}{#1 --- #2}}%
}

\hypersetup{
  colorlinks=true,
  linkcolor=black,
  urlcolor=blue!60!black,
  citecolor=blue!60!black,
  pdftitle={Ziren Multilinear PCS Pipeline: From WHIR Soundness Fixes to a BaseFold Production Backend},
  pdfauthor={ProjectZKM Research}
}

\definecolor{codebg}{gray}{0.95}
\definecolor{codeborder}{gray}{0.85}

\lstset{
  basicstyle=\ttfamily\footnotesize,
  backgroundcolor=\color{codebg},
  frame=single,
  rulecolor=\color{codeborder},
  framesep=4pt,
  xleftmargin=6pt,
  xrightmargin=6pt,
  breaklines=true,
  breakatwhitespace=false,
  postbreak=\mbox{\textcolor{gray}{$\hookrightarrow$}\space},
  columns=fullflexible,
  keepspaces=true,
  aboveskip=0.6em,
  belowskip=0.6em,
  showstringspaces=false,
}

\renewcommand{\abstractname}{\large Abstract}

\title{\Large\textbf{Ziren Multilinear PCS Pipeline}\\[0.3em]
  \large From WHIR Soundness Fixes to a BaseFold Production Backend%
  \\[0.5em]
  \normalsize with Jagged Packing, Late-Binding, and Per-Chip Commitments}

\author[1]{ProjectZKM Research}
\affil[1]{ProjectZKM}

\date{April 2026}

\begin{document}
\maketitle

\begin{abstract}
\noindent
We document a two-phase upgrade of the Ziren zkVM's multilinear
polynomial commitment stack.  The first phase is a production-sound
integration of the WHIR multilinear PCS, closing three soundness
gaps in the upstream WHIR fast path: (i)~unbound transition
constraints, fixed by a sumcheck-based zerocheck; (ii)~unbound
lookup arguments, fixed by a sumcheck reduction of LogUp-GKR; and
(iii)~unbound per-row MLE openings, fixed by a \emph{late-binding}
WHIR adapter that allows evaluation points to be sampled
\emph{after} commitment.  A \emph{cross-binding} spot-check
construction ties the FRI commit and the WHIR late-binding commit
to the same trace, ruling out a divergent-trace attack.  A
\emph{jagged late-binding} reduction combines variable-height
chip-trace packing with the late-binding adapter, achieving
$57$--$85\times$ smaller per-shard proofs at $1.3$--$2.1\times$
prover cost.

The second phase replaces WHIR entirely with a
\emph{BaseFold}-backed stack (E1 WHIR removal, April 2026).  The
BaseFold port matches SP1's \texttt{slop\_basefold} upstream;
jagged packing stays as the heterogeneous-batch adapter.  The new
per-chip variant (E3, in-progress) avoids the dense materialisation
the WHIR path required, eliminating a structural out-of-memory
failure mode and replacing a 2000\,LOC reduction with 16 focused
unit tests validating each primitive (LongMle, HadamardProduct,
sumcheck driver, BranchingProgram, and end-to-end single-chip +
multi-chip jagged-eval protocol runs).

All host-side code is implemented and tested in the open-source
Ziren codebase: the late-binding modules ($24$ WHIR-era tests
preserved as reference, passing), the recursion-circuit verifier
($18$ tests), zerocheck + LogUp-GKR ($9$ tests), and the BaseFold
stack ($56/56$ in \texttt{zkm-stark} including $16$ in
\texttt{jagged\_per\_chip}).  We document the workload-dependent
trade-off from WHIR profiling, motivate the migration, and report
BaseFold performance (fib-1k: $9.6$\,s prove, $362$\,ms verify,
$4.96$\,MB; keccak-precompile SDK: $50.9$\,s wall, $10.3$\,GB peak,
verified end-to-end).

\paragraph{Keywords:}
STARKs $\cdot$ multilinear polynomial commitment $\cdot$ WHIR
$\cdot$ BaseFold $\cdot$ jagged commitment $\cdot$ soundness $\cdot$
recursion $\cdot$ zero-knowledge virtual machines $\cdot$ Plonky3.
\end{abstract}

\tableofcontents
\vspace{1em}
\hrule
\vspace{1em}

% =============================================================================
% Section 1: Introduction
\section{Introduction}
\label{sec:introduction}

Zero-knowledge virtual machines (zkVMs) compile general-purpose
computation into succinct cryptographic proofs.  The Ziren
zkVM~\cite{ziren} is a MIPS-based zkVM built on the
Plonky3~\cite{plonky3} STARK toolkit, currently shipping with FRI as
its polynomial commitment scheme (PCS).  Although FRI offers
excellent verification time, modern multilinear PCS such as
WHIR~\cite{whir2024} promise simultaneously smaller proofs,
faster verification, and a more natural multilinear-extension (MLE)
view of trace data --- all of which are increasingly relevant as
zkVMs adopt hypercube-style architectures~\cite{sp1-hypercube}.

This paper documents the cryptographic and engineering work required
to migrate Ziren's STARK pipeline from FRI to a production-sound
WHIR-based pipeline.  The transition is non-trivial: the existing
``WHIR fast path'' shipped in Ziren is a structural skip of
permutation traces and quotient evaluation, with several unbound
soundness gaps.  We close all of them.

\subsection{Contributions}

\begin{enumerate}
  \item \textbf{Sumcheck-based zerocheck for transition constraints
        (\S\ref{sec:zerocheck-phase2a}).}  Replaces the quotient
        polynomial argument with a sumcheck protocol
        proving $\sum_{b \in \{0,1\}^m} \mathsf{eq}(r, b) \cdot
        C(b) = 0$, with provable soundness $O(d \cdot m / |\mathbb{F}|)$
        for constraint degree $d$.

  \item \textbf{Sumcheck-based LogUp-GKR for lookup
        arguments (\S\ref{sec:logup-gkr-fix}).}  We identify and
        fix a critical soundness bug in the initial folding-only
        implementation: the combine identity
        $(N(z\Vert 0)\cdot D(z\Vert 1) + D(z\Vert 0) \cdot
         N(z\Vert 1), D(z\Vert 0) \cdot D(z\Vert 1))$
        is degree-$2$ in each coordinate, which only matches the
        verifier's degree-$1$ multilinear $N_{k+1}, D_{k+1}$ at
        Boolean points.  Our replacement protocol uses a per-layer
        sumcheck reduction over a degree-$3$ round polynomial to
        ensure soundness at random points.

  \item \textbf{Late-binding WHIR adapter
        (\S\ref{sec:whir-late-binding}).}  The upstream WHIR
        \texttt{MultilinearPcs} adapter requires evaluation points
        \emph{at commit time}, which is incompatible with our
        protocol where the per-chip row coordinates
        $r_{\mathrm{row}}$ are derived \emph{after} commit via GKR.
        We construct an in-tree adapter built directly on
        \texttt{p3-whir} internals (\texttt{CommitmentWriter},
        \texttt{Verifier}, \texttt{EqStatement}) that exposes
        single-column and multi-column commit + post-commit
        \texttt{add\_late\_claim} + open API surfaces.  The
        cryptographic substrate is unchanged from upstream WHIR;
        only the transcript ordering differs.

  \item \textbf{Cross-binding (C.1) Brakedown spot-check
        protocol.}  In a dual-commit setting where FRI and WHIR
        both commit to the same trace, a malicious prover can
        commit divergent traces $A \neq B$.  We close this with a
        Brakedown-style spot-check argument~\cite{brakedown}: $K$
        random trace positions are sampled post-commit; the prover
        provides merkle paths against \emph{both} commits at each
        position.  Tampering on either side is detected with
        soundness $1 - (1 - 1/N)^K$.  For $K = 80$ and $N = 2^{17}$
        this gives $\sim 2^{-80}$ missed-detection probability.

  \item \textbf{Jagged late-binding construction.}  We combine
        Jagged~\cite{jagged} variable-height trace packing with the
        late-binding adapter via a sumcheck reduction that ties
        per-chip per-column row-MLE claims to a single dense
        commitment.  This collapses $\sim 570$ per-chip-per-column
        WHIR commits to a single shard-level commit, reducing
        proof size by a factor of $57$--$85\times$ in our
        evaluation (\S\ref{sec:performance}).

  \item \textbf{Recursion-circuit WHIR verifier scaffolding.}  We
        port the host-side verification logic into the
        recursion-circuit through a series of focused sub-systems:
        a \texttt{SumcheckVerify} chip (full DSL$\to$AIR pipeline,
        seven layers), a merkle-path verifier built atop the
        existing Poseidon2 chip, OOD evaluation primitives, and
        STIR query position derivation.  Each sub-system is
        independently testable and serves as a stepping stone
        toward in-circuit \texttt{verify\_whir\_pcs}.
\end{enumerate}

\subsection{Implementation and Open-Source Release}

All host-side prover and verifier code is implemented in the open-
source Ziren codebase, opt-in via two environment variables:
\texttt{ZIREN\_USE\_WHIR=1} (selects the WHIR fast path with
zerocheck and LogUp-GKR), and \texttt{ZIREN\_LATE\_BINDING=jagged}
(selects the jagged late-binding path).  A third flag
\texttt{ZIREN\_CROSS\_BINDING=1} activates the cross-binding spot-
check protocol.  Defaults preserve the original FRI pipeline so
existing deployments are unaffected.

We test each cryptographic mechanism via end-to-end round-trip and
tamper-rejection unit tests, plus end-to-end smoke tests on the
fibonacci, JSON, and keccak workloads.  At the time of writing,
$24$ late-binding-related tests, $9$ zerocheck and LogUp-GKR tests,
$18$ recursion-circuit WHIR verifier tests, and the existing $34$
recursion-core tests all pass.  Compressed-proof verification of
each smoke-tested workload completes successfully.

\subsection{Paper Outline}

\S\ref{sec:preliminaries} reviews the necessary preliminaries
(STARK, sumcheck, WHIR, LogUp-GKR, Plonky3 architecture).
\S\ref{sec:background} sets the engineering context, summarises
the Plonky3 upgrade work, and tabulates implementation status.
\S\ref{sec:plonky3-upgrade} details the recursion-circuit FRI
verifier upgrade required by upstream Plonky3 protocol changes.
\S\ref{sec:soundness-fix} formalises the threat model of the
existing WHIR fast path.
\S\S\ref{sec:zerocheck-phase2a}--\ref{sec:whir-late-binding}
present the three new constructions: zerocheck (Phase 2a),
LogUp-GKR (Phase 2b), and late-binding WHIR with cross-binding
and jagged packing (Phase 2c+).
\S\ref{sec:performance} reports performance measurements.
\S\ref{sec:related-work} positions the work against prior
constructions.  \S\ref{sec:conclusion} concludes with future
directions.  Appendices document detailed implementation status
and a session-by-session changelog.


% =============================================================================
% Section 2: Preliminaries
\section{Preliminaries}
\label{sec:preliminaries}

\subsection{Notation}

We work over a finite field $\mathbb{F}$ (KoalaBear in our deployment,
$|\mathbb{F}| \approx 2^{31}$) and an extension $\mathbb{F}_q$ with
$q = |\mathbb{F}|^D$ for $D=4$, giving challenger soundness
$\sim 2^{-124}$.  $N = 2^n$ denotes a power-of-two domain size,
typically the height of a chip's trace.  Bold lowercase letters
($\mathbf{x}, \mathbf{r}, \mathbf{z}$) denote vectors of field
elements; uppercase ($P, Q, F$) denote polynomials.

For a vector of evaluations $\mathbf{v} \in \mathbb{F}^{2^m}$, the
\emph{multilinear extension} (MLE) $\widetilde{\mathbf{v}}:
\mathbb{F}^m \to \mathbb{F}$ is the unique multilinear polynomial
satisfying $\widetilde{\mathbf{v}}(\mathbf{b}) = v_{\mathrm{idx}
(\mathbf{b})}$ for $\mathbf{b} \in \{0,1\}^m$.  We use the
multilinear equality polynomial $\mathsf{eq}(\mathbf{r}, \mathbf{x})
= \prod_{i=1}^{m} (r_i x_i + (1-r_i)(1-x_i))$, which satisfies
$\widetilde{\mathbf{v}}(\mathbf{r}) = \sum_{\mathbf{b} \in \{0,1\}^m}
\mathsf{eq}(\mathbf{r}, \mathbf{b}) \cdot v_{\mathrm{idx}(\mathbf{b})}$.

\subsection{STARKs and the Ziren Pipeline}

A STARK proves correct execution of a virtual-machine program by:
(1) expressing the program's per-cycle behaviour as an algebraic
intermediate representation (AIR), (2) committing to the resulting
trace via a polynomial commitment scheme (PCS), (3) verifying that
the committed trace satisfies the AIR's transition and boundary
constraints (via a quotient polynomial argument or a sumcheck),
and (4) verifying lookups (via the LogUp~\cite{logup-gkr} family of
arguments).

Ziren's main branch implements step~(2) via FRI on a multiplicative
subgroup, with separate commits to main, permutation, and quotient
traces.  Step~(3) is the standard quotient identity
$Q(\zeta) \cdot Z_H(\zeta) = C_{\mathrm{batched}}(\zeta)$ at a
random univariate point $\zeta$.

\subsection{Sumcheck Protocol}

The sumcheck protocol~\cite{sumcheck} reduces a claim about a
polynomial sum over the Boolean hypercube to a single polynomial
evaluation:
\[
  \text{Claim:}\quad \sum_{\mathbf{b} \in \{0,1\}^m} F(\mathbf{b}) = v.
\]
For each round $i \in [m]$, the prover sends a univariate polynomial
$p_i(X) = \sum_{\mathbf{b}' \in \{0,1\}^{m-i}}
F(r_1, \ldots, r_{i-1}, X, \mathbf{b}')$ of degree $\deg(F)$, and
the verifier checks $p_i(0) + p_i(1) = s_i$ (where $s_1 = v$ and
$s_{i+1} = p_i(r_i)$ for verifier-sampled $r_i \in \mathbb{F}_q$).
After $m$ rounds, the verifier checks
$F(r_1, \ldots, r_m) = s_{m+1}$.  Soundness error is
$d \cdot m / |\mathbb{F}_q|$ for $\deg(F) = d$.

\subsection{WHIR}

WHIR~\cite{whir2024} is an interactive oracle proof of proximity
(IOPP) for Reed--Solomon (RS) codes that supports multilinear
polynomial commitments with $O(\log n)$ verifier work and proof
sizes $5$--$10\times$ smaller than FRI at equivalent security.

A WHIR commitment to $f \in \mathbb{F}^{2^n}$ encodes $f$'s
coefficients into an RS codeword on a coset of size $2^{n+r}$
($r$ = log-inv-rate) and merkle-commits the row hashes.  The open
phase combines sumcheck folding rounds with STIR query openings.

The Plonky3 implementation \texttt{p3-whir} provides a
\texttt{MultilinearPcs} adapter requiring opening points
\emph{at commit time} (absorbed into Fiat-Shamir before the
batching challenge $\alpha$ is sampled).  This is the standard
soundness ordering when opening points are pre-known; our
work (\S\ref{sec:whir-late-binding}) constructs an alternative
\emph{late-binding} ordering for protocols where the opening
point is derived post-commit.

\subsection{LogUp-GKR}

LogUp~\cite{logup-gkr} reduces a multiset-equality lookup argument
to a sum-of-fractions identity
\[
  \sum_{i \in \mathrm{senders}} \frac{m_i}{\alpha - f_i} \;=\;
  \sum_{j \in \mathrm{receivers}} \frac{m_j}{\alpha - f_j},
\]
where $m_*$ are multiplicities and $f_*$ are interaction
fingerprints.  GKR~\cite{logup-gkr} verifies this fraction-sum
identity via a layered tree of fraction combinations
$(a, b) \oplus (c, d) = (ad + bc, bd)$, opening only the root
$(\mathrm{num}, \mathrm{denom})$ and verifying $\mathrm{num} = 0$.

The verification reduces tree-root claims to leaf claims via
sumcheck rounds at each layer.  The leaf claims are MLE evaluations
of multiplicities and fingerprints at a sumcheck-derived point
$r_{\mathrm{row}}$, which is therefore \emph{post-commit}.  This
post-commit derivation is the central structural difficulty for
binding the leaf claims to a WHIR commitment of the trace.

\subsection{Plonky3 Architecture}

Ziren is built on Plonky3~\cite{plonky3}, which provides chip-based
AIR composition, Fiat-Shamir challenger, MMCS (mixed-matrix
commitment scheme) merkle trees, and a recursion compiler.  Our
work modifies neither the chip composition nor the recursion
compiler core; we add new chips (sumcheck verification),
new instructions (sumcheck verify), new prover/verifier paths
(WHIR fast path), and new helpers (late-binding, cross-binding,
jagged + late-binding) on top of the existing infrastructure.


% =============================================================================
% Section 3: Background and Motivation (status / context)
\section{Background and Motivation}
\label{sec:background}

\subsection{The starting point}

Ziren's main branch ships a STARK prover with FRI as the polynomial
commitment scheme.  Each shard commits the main trace, the permutation
trace (for LogUp lookups), and a quotient polynomial (for transition
constraints), then opens all three at a univariate point $\zeta$.  The
verifier replays the same three opens and checks the AIR identities:
constraint-quotient relationship, lookup accumulation, cumulative-sum
closure.

This is the standard univariate-FRI STARK pipeline.  It has the
disadvantages typical of FRI: large per-shard proofs (megabytes for
RISC-V workloads), expensive recursion (many merkle paths per query),
and a large recursion circuit (the FRI verifier must be expressed in
the recursion AIR, with all the per-round PoW witnesses, log-arity
observations, and merkle-path checks).

\subsection{The Plonky3 / WHIR upgrade}

Two upstream developments motivated a major upgrade:

\begin{enumerate}
  \item \textbf{Plonky3 protocol updates.}  Upstream Plonky3 added
        per-round PoW witnesses, log-arity observations, and a
        restructured opened-values absorption pattern to the FRI
        verifier.  This is necessary for proper Fiat-Shamir binding
        with adaptive folding factors but breaks compatibility with
        any recursion circuit pinned to the older protocol.
  \item \textbf{WHIR multilinear PCS.}  The WHIR protocol
        \cite{whir2024} reduces multilinear polynomial commitments to
        a sequence of sumcheck rounds with folding, achieving
        $O(\log n)$ verifier work and proof sizes 5--10$\times$
        smaller than FRI for the same security level.  Plonky3's
        \texttt{p3-whir} adapter exposes WHIR as a
        \texttt{MultilinearPcs}.
\end{enumerate}

\subsection{The architecture of this work}

The work documented in this whitepaper does three things, in order:

\begin{enumerate}
  \item \textbf{Plonky3 upgrade.}  Bring Ziren's recursion circuit and
        FRI verifier in line with the new Plonky3 protocol.  Includes
        per-round PoW witnesses in the recursion FRI verifier, fixed
        \texttt{check\_witness} for empty PoW, and a \texttt{is\_div\_active}
        preprocessed column to gate \texttt{DivF}/\texttt{DivE} in the
        recursion AIR (see \S\ref{sec:plonky3-upgrade}).
  \item \textbf{Phase 2 soundness fix.}  Replace the quotient and
        permutation-trace mechanisms with sumcheck-based equivalents
        that work multilinearly: zerocheck (Phase 2a) and LogUp-GKR
        (Phase 2b).  This is what makes the WHIR fast path actually
        sound, not just a structural skip
        (see \S\ref{sec:zerocheck-phase2a}, \S\ref{sec:logup-gkr-fix}).
  \item \textbf{Late-binding WHIR adapter (Phase 2c).}  The remaining
        soundness gap after Phase 2b: per-chip row-MLE openings
        (\texttt{LogUpRowOpening}) are sent in the proof but not
        cryptographically tied to the main-trace commitment.  Closing
        this gap requires opening the WHIR commitment at multilinear
        points sampled \emph{after} the commit, which the upstream
        adapter's API forbids.  This work introduces a late-binding
        adapter that correctly handles post-commit point selection
        (see \S\ref{sec:whir-late-binding}).
\end{enumerate}

\subsection{Status of the work}

\begin{tabularx}{\textwidth}{@{}lXc@{}}
\toprule
\textbf{Item} & \textbf{Status} & \textbf{Tests} \\
\midrule
Plonky3 upgrade (recursion FRI verifier) & shipped, compressed proofs verify & --- \\
Cherry-pick aggregation fix (\texttt{7c2071e}) & shipped on \texttt{feat/upgrade-plonky3} & --- \\
Aggregation example with WHIR (ZK stages) & all stages succeed; gnark wrap blocked on stale artifacts & --- \\
Groth16 example with WHIR (ZK stages) & all stages succeed; gnark wrap blocked on stale artifacts & --- \\
Phase 2a (zerocheck) & shipped & 5/5 passing \\
Phase 2b (LogUp-GKR sumcheck reduction) & shipped & 9/9 passing \\
Phase 2c (single-column late-binding adapter) & shipped & 2/2 passing \\
Phase 2c (multi-column late-binding adapter) & shipped & 2/2 passing \\
Phase 2c (\texttt{LateBindingCapable} trait + KB impl) & shipped & --- \\
Phase 2c (full multi-chip wiring simulation) & shipped & 1/1 passing \\
Phase 2c (JSON serialisation + round-trip test) & shipped & 1/1 passing \\
Phase 2c (interleaved-codeword optimisation, VEIL Lemma 4.7) & designed; implementation deferred & --- \\
Prover wiring (\texttt{prover.rs} \texttt{try\_compute\_late\_binding\_proofs}) & shipped (TypeId dispatch, dual commit, dispatches both \texttt{KoalaBearPoseidon2} and \texttt{...Inner}) & --- \\
Verifier wiring (\texttt{verifier.rs} \texttt{try\_verify\_late\_binding\_proofs}) & shipped (TypeId dispatch, mirroring fresh challenger) & --- \\
End-to-end smoke test (fibonacci with \texttt{ZIREN\_USE\_WHIR=1}) & verified: 19 chips, 29.5\,MB late-binding bytes, 32.6\,s prover, verifier accepts & --- \\
Cross-binding (FRI $\leftrightarrow$ WHIR consistency) & not started; soundness gap documented & --- \\
Recursion circuit (\texttt{verify\_whir\_pcs}) & not started & --- \\
gnark BN254 wrap artifact regeneration & blocked on VK map regen & --- \\
\bottomrule
\end{tabularx}


% =============================================================================
% Section 4: Plonky3 Upgrade
\section{Plonky3 Upgrade}
\label{sec:plonky3-upgrade}

The recursion circuit needed several fixes to work with the new
upstream FRI verifier.  All fixes are committed on
\texttt{feat/upgrade-plonky3}.

\subsection{Per-round PoW witnesses}

Upstream's new FRI requires the prover to grind a proof-of-work
witness once per round (not just at the final query stage).  The
recursion-circuit FRI verifier
(\texttt{crates/recursion/circuit/src/fri.rs}) was missing the
per-round \texttt{check\_witness} call.  Added.

\subsection{Log-arity observations and opened-values restoration}

The new FRI verifier observes the log-arity of each round and absorbs
the opened-values vector into the transcript before sampling the next
fold randomness.  Both observations were stripped in an earlier
"compatibility" patch and were restored in
\texttt{p3-whir} commit \texttt{715fdb5}.  The recursion circuit
mirrors them.

\subsection{\texttt{check\_witness} early-return for zero PoW bits}

\texttt{GrindingChallenger::check\_witness} returns \texttt{true} when
the requested bits is zero.  The recursion circuit's variable-arity
version was running the hash anyway; aligned to native behaviour.

\subsection{\texttt{is\_div\_active} preprocessed column}

The recursion compiler emits \texttt{DivF}/\texttt{DivE} instructions
in both branches of \texttt{Select}, including dead branches where the
divisor is zero.  Native execution handles this gracefully because
\texttt{mult==0} skips the operation.  The AIR didn't.  Added a
preprocessed column \texttt{is\_div\_active = is\_div AND
mult\_active} (where \texttt{mult\_active} marks active alu rows) and
gated the division constraint on it.  Mirror runtime guard added in
\texttt{crates/recursion/core/src/runtime/mod.rs}.

\subsection{Result}

Compressed proofs for fibonacci and other benchmarks verify with the
upgraded recursion circuit.  Aggregation runs through
\texttt{prove\_core $\to$ compress $\to$ shrink $\to$ wrap\_bn254}
successfully; the gnark BN254 wrap fails with witness-count mismatch
because the gnark circuit artifacts predate the FRI changes (250
extra witness elements; see \S\ref{sec:performance}).


% =============================================================================
% Section 5: WHIR Fast Path Threat Model
\section{Soundness Fix Design: Making WHIR Mode Production-Sound}
\label{sec:soundness-fix}

\subsection{Current State}

The WHIR fast path (\texttt{ZIREN\_USE\_WHIR=1}) in
\texttt{crates/stark/src/prover.rs} skips transition-constraint verification
entirely. The core prover:
\begin{enumerate}
  \item Commits to the main trace only (no permutation trace, no quotient).
  \item Opens the main trace at $\zeta$ via WHIR (Jagged).
  \item Sets \texttt{ShardCommitment.permutation\_commit = None} and
        \texttt{quotient\_commit = None}.
  \item Leaves \texttt{ChipOpenedValues.permutation} and \texttt{quotient}
        empty.
\end{enumerate}

The recursion verifier (both native and circuit) short-circuits when
\texttt{permutation\_commit.is\_none()}: it skips the PCS round for
permutation/quotient, skips \texttt{verify\_constraints}, and only checks:
\begin{itemize}
  \item PCS opening of preprocessed + main.
  \item $\sum_i \text{local\_cumulative\_sum}_i = 0$.
\end{itemize}

\subsection{What This Leaves Unchecked}

Given a committed main trace $M$ and cumulative-sum extraction from the last
14 base-field columns of $M$, the verifier does \emph{not} check:

\begin{table}[h]
\centering
\small
\begin{tabular}{lcc}
\toprule
\textbf{AIR check} & \textbf{FRI mode} & \textbf{WHIR fast path} \\
\midrule
$\sum \text{lookup events} = 0$ (permutation argument) &
  \checkmark\ via quotient check &
  \checkmark\ via cumulative sum \\
Transition: $C(M[i], M[i+1], P[i]) = 0 \ \forall i$ &
  \checkmark\ via quotient check &
  $\times$\ NOT CHECKED \\
Boundary: $C_{\text{first}}(M[0]) = 0$, $C_{\text{last}}(M[n{-}1]) = 0$ &
  \checkmark\ via Lagrange selectors &
  $\times$\ NOT CHECKED \\
Public values = $M[0][\text{pub\_cols}]$ &
  \checkmark\ via boundary constraint &
  $\times$\ NOT CHECKED \\
\bottomrule
\end{tabular}
\caption{What is verified in each mode.}
\label{tab:soundness-gap}
\end{table}

A malicious prover can therefore produce a ``verifying'' proof for a main
trace that does not come from any valid execution of the MIPS program,
forge public outputs, or violate transition constraints. Only the global
septic digest (extracted from the last 14 columns) is bound, and that is
bindable separately to memory-init/finalize events.

\emph{The current WHIR mode is a benchmark prototype, not a zk-proof.}

\subsection{Target Soundness Design}

The fix must re-introduce transition-constraint verification while
preserving the Jagged + WHIR performance characteristics. There are three
standard options.

\paragraph{Option A --- Keep the quotient polynomial, batch it into WHIR.}
Commit to the quotient polynomial alongside the main trace, open both at
$\zeta$, check $q(\zeta) \cdot Z_H(\zeta) = \text{folded\_constraints}(\zeta)$.
\emph{Pro:} minimal change. \emph{Con:} defeats most of the WHIR speedup ---
quotient width is $2^{\text{log\_quotient\_degree}}$ (= 2 in Ziren), each
column must be computed, committed, and opened. The tendermint $4.5\times$
speedup would collapse to $\sim$1.3$\times$. \emph{Status:} not recommended.

\paragraph{Option B --- Logup-GKR plus sumcheck-based constraint verification
(SP1 hypercube model).}
Replace both the permutation argument and the quotient polynomial with GKR /
sumcheck reductions:
\begin{enumerate}
  \item \textbf{Constraints $\to$ sumcheck.} The transition-constraint
        polynomial $C(X) = \alpha \cdot C_0(X) + \alpha^2 \cdot C_1(X) + \cdots$
        is viewed as a multilinear polynomial over the Boolean hypercube
        $\{0,1\}^{\log n}$. The prover proves
        $$
          \sum_{x \in \{0,1\}^{\log n}} \mathsf{eq}(r, x) \cdot C(x, x_{\text{next}}) = 0
        $$
        via sumcheck ($\log n$ rounds, $\sim$3 field elements per round).
  \item \textbf{Lookups $\to$ Logup-GKR.} The lookup identity becomes a GKR
        circuit whose final claim is a multilinear evaluation of a fraction
        polynomial. The prover runs a few sumcheck rounds to reduce it to a
        single evaluation.
  \item \textbf{Final opening.} All final claims are batched and opened at
        one multilinear point via WHIR.
\end{enumerate}
\emph{Pro:} matches SP1's direction; Logup-GKR has battle-tested designs;
preserves most of the WHIR speedup because no extra trace is committed.
\emph{Con:} large implementation effort (requires sumcheck-verifier chips
inside the recursion circuit). \emph{Status:} recommended (this document's
focus).

\paragraph{Option C --- Multilinear constraint check via WHIR opening at $\zeta$.}
Keep the univariate STARK constraint check but express it in multilinear
form. The verifier:
\begin{enumerate}
  \item Samples $\alpha, \zeta$.
  \item Opens $\text{main}(\zeta), \text{main}(\zeta \omega), \text{preproc}(\zeta), \text{preproc}(\zeta \omega)$ via WHIR.
  \item Locally computes $\text{folded\_constraints}(\zeta)$ from the opened values.
  \item Asserts $\text{folded\_constraints}(\zeta) = 0$.
\end{enumerate}
The last step replaces the identity $q(\zeta) \cdot Z_H(\zeta) = \text{folded}(\zeta)$
with the direct assertion $\text{folded}(\zeta) = 0$. This is only sound if
$\text{folded\_constraints}(X)$ vanishes on the trace domain, which still
requires a sumcheck argument to prove --- collapsing Option C into
Option B. \emph{Status:} subsumed by Option B.

\subsection{Chosen Design: Option B}

\subsubsection*{Protocol sketch (prover)}

Let $M$ be the main trace (multilinear, dimension $\log n$ over
$\{0,1\}^{\log n}$), $P$ the preprocessed trace, $\mathrm{pub}$ the public
values.

\begin{enumerate}[label=\textbf{Round \arabic*.}, wide=0pt]
  \item Commit to $M$ via WHIR (same as today).
  \item Sample constraint-mixing challenge $\alpha$.
  \item \textbf{Constraint sumcheck.} Prover and verifier run sumcheck for
        $$
          \sum_{x \in \{0,1\}^{\log n}} \mathsf{eq}(r, x) \cdot \Bigl[ \sum_i \alpha^i \cdot C_i\bigl(M(x), M(x \omega), P(x), P(x \omega), \mathrm{pub}\bigr) \Bigr] = 0,
        $$
        where $r \in \mathbb{F}^{\log n}$ is sampled by the verifier at the
        start. After $\log n$ rounds, the claim reduces to a single
        multilinear-value claim on $M$ and $P$ at point $r' \in \mathbb{F}^{\log n}$.
  \item \textbf{Lookup (Logup-GKR).} For each lookup table $T$ with events
        $e_1, \dots, e_k$:
        $$
          \sum_j \frac{1}{\beta - e_j} = \sum_t \frac{\text{multiplicity}(t)}{\beta - t}.
        $$
        The prover runs a GKR circuit whose output wire is $0$. The GKR
        reduction ends with a claim on $M, P$ at another point $r''$.
  \item \textbf{Batched WHIR opening.} All multilinear claims on $M$ (from
        sumchecks and GKR) are batched via random linear combination and
        opened at a single point via WHIR. Same for $P$.
\end{enumerate}

\subsubsection*{Verifier cost}

\begin{itemize}
  \item Sumcheck of degree-$d$ polynomial over $\log n$ rounds:
        $\log n \cdot (d+1)$ field elements observed, one final multilinear
        evaluation check.
  \item Logup-GKR: $O(\log k)$ sumcheck rounds per table.
  \item WHIR opening: same cost as today's Jagged + WHIR open.
\end{itemize}
Total verifier work: $O\bigl(\log n \cdot (d + |\text{tables}|)\bigr)$ ---
roughly the same asymptotic cost as today, \emph{plus} constraint
soundness.

\subsubsection*{Recursion implications}

The recursion circuit currently contains:
\begin{itemize}
  \item \texttt{verify\_shard} --- iterates over chips, calls
        \texttt{verify\_constraints}.
  \item \texttt{verify\_constraints} --- computes
        $\text{folded\_constraints}(\zeta)$ and
        $q(\zeta) \cdot Z_H(\zeta)$ and asserts equality.
\end{itemize}

Under Option B the recursion circuit needs:
\begin{enumerate}
  \item A \texttt{SumcheckVerify} chip (skeleton already exists in
        \texttt{crates/recursion/core/src/chips/sumcheck\_verify.rs}).
  \item A \texttt{GKRVerifier} circuit (new).
  \item A \texttt{MultilinearEvalCheck} that reconciles a claimed
        multilinear evaluation at $r$ with the value opened by WHIR.
  \item A replacement for \texttt{verify\_constraints} that runs
        \texttt{SumcheckVerify} per chip, runs \texttt{GKRVerifier} per
        lookup table, batches all multilinear claims, and checks them
        against WHIR-opened values.
\end{enumerate}

\subsubsection*{Migration path}

\begin{enumerate}
  \item \textbf{Phase 2a (done)} --- sumcheck primitives, hypercube
        constraint evaluator, \texttt{ShardProof::zerocheck\_proofs} field,
        prover + verifier wiring. See Section~\ref{sec:zerocheck-phase2a}.
  \item \textbf{Phase 2b} --- Logup-GKR for lookup arguments in the native
        prover. Removes the permutation trace entirely.
  \item \textbf{Phase 2c} --- wire \texttt{SumcheckVerify} and
        \texttt{GKRVerifier} into the recursion circuit. Replace
        \texttt{verify\_constraints} with the sumcheck-based version.
  \item \textbf{Phase 3} --- formal verification of the sumcheck / GKR
        verifier chips.
\end{enumerate}

\subsection{Projected Performance Impact}

Compared to the current (unsound) WHIR fast path:

\begin{table}[h]
\centering
\small
\begin{tabular}{lllc}
\toprule
\textbf{Operation} & \textbf{WHIR fast path} & \textbf{WHIR + Logup-GKR} & \textbf{Overhead} \\
\midrule
Main trace commit   & same   & same    & 0\%       \\
Constraint check    & skipped & sumcheck ($\log n$ rounds) & +10--15\% \\
Lookup check        & cumulative sum read & Logup-GKR sumcheck & +5--10\%  \\
PCS open            & one point & batched multi-point & +5\%      \\
\midrule
\textbf{Total expected} & \textbf{baseline} & \multicolumn{2}{r}{\textbf{+20--30\%}} \\
\bottomrule
\end{tabular}
\caption{Projected end-to-end overhead of Option B relative to the unsound
WHIR fast path.}
\label{tab:sound-overhead}
\end{table}

Projected tendermint performance: $502\,\text{s} \cdot 1.25 \approx 630\,\text{s}$,
still $\mathbf{3.6\times}$ faster than main FRI (2270\,s), with full
soundness restored.

\subsection{Immediate Action Items}
\begin{enumerate}
  \item Rename the current flag to \texttt{ZIREN\_WHIR\_UNSAFE=1} with a
        runtime warning; the existing name misleads about soundness.
  \item Complete Phase 2a follow-up: multi-point WHIR opening so the
        sumcheck final claim $C(r')$ can be checked against an actual
        opening.
  \item Implement Phase 2b (Logup-GKR) to make all chips zerocheck-eligible.
  \item Wire into the recursion circuit (Phase 2c).
  \item Formal-verification pass (Phase 3).
\end{enumerate}


% =============================================================================
% Section 6: System Design --- proof-system architecture (jagged + WHIR + late-
% binding + cross-binding composition, with implementation pointers)
\section{System Design: Proof-System Architecture}
\label{sec:design}

This section presents the full proof-system design of the Ziren WHIR
pipeline.  Where prior sections (\S\S\ref{sec:zerocheck-phase2a}--\ref{sec:whir-late-binding})
present each construction in isolation, here we describe how they
compose into a single end-to-end prover/verifier flow, with explicit
references to the implementation files.

\subsection{Proof-System Overview}

A Ziren shard proof attests that one shard of a MIPS program
executed correctly.  The proof contains:
\begin{itemize}
  \item A polynomial commitment to the per-chip main traces.
  \item Sumcheck-based proofs of transition-constraint and lookup
        soundness.
  \item Per-row multilinear-extension (MLE) openings of the main
        trace at GKR-derived points, bound back to the main commit
        via a late-binding wrapper.
  \item Cross-binding spot-check evidence linking the FRI commit
        and the WHIR late-binding commit to the same trace
        (production target).
\end{itemize}

The pipeline replaces three FRI-pipeline components with multilinear
equivalents:
\begin{enumerate}
  \item \textbf{Quotient identity} ($Q(\zeta) \cdot Z_H(\zeta) =
        C(\zeta)$) is replaced by zerocheck:
        $\sum_{\mathbf{b} \in \{0,1\}^m} \mathsf{eq}(\mathbf{r},
        \mathbf{b}) \cdot C(\mathbf{b}) = 0$, verified via sumcheck.
  \item \textbf{Permutation-trace lookup arguments} are replaced by
        LogUp-GKR with a per-layer sumcheck reduction.
  \item \textbf{Univariate FRI opening at $\zeta$} is augmented (not
        replaced --- the dual-commit approach keeps both for now)
        by WHIR multilinear opening at the GKR-derived row coordinates
        $\mathbf{r}_{\mathrm{row}}$.
\end{enumerate}

\subsection{Data Layout: Jagged Polynomial Commitment}
\label{sec:design-jagged}

Each shard contains $N$ chip traces of differing heights $h_c$ and
widths $w_c$, $c = 0, \ldots, N-1$.  Naively committing each chip
separately requires $N$ merkle trees and $N$ PCS instances.  We
adopt the Jagged construction~\cite{jagged} to amortise this:

\begin{construction}[Jagged trace packing]
\label{cons:jagged-pack}
Let $\mathbf{T}_0, \ldots, \mathbf{T}_{N-1}$ be per-chip trace
matrices.  The packed dense vector $\mathbf{q}$ is the column-major
concatenation of all per-chip columns:
\[
  \mathbf{q} \;=\;
  \mathbf{T}_0[:,0] \,\Vert\, \mathbf{T}_0[:,1] \,\Vert\, \ldots
   \,\Vert\, \mathbf{T}_{N-1}[:,w_{N-1}-1].
\]
Cumulative offsets $t_k$ satisfy $t_0 = 0$, $t_k = t_{k-1} +
h_{c(k)}$ where $k$ ranges over global column indices.  Padding
brings $|\mathbf{q}|$ to the next power of two $2^n$.
\end{construction}

The sparse $(\mathrm{row}, \mathrm{col})$-indexed polynomial $p(\mathbf{z}_r,
\mathbf{z}_c)$ relates to $\mathbf{q}$'s MLE $\widetilde{\mathbf{q}}$ via
\[
  p(\mathbf{z}_r, \mathbf{z}_c) \;=\;
  \sum_{j=0}^{|\mathbf{q}|-1}
    \widetilde{\mathbf{q}}(j) \cdot
    \mathsf{eq}(\mathrm{row}(j), \mathbf{z}_r) \cdot
    \mathsf{eq}(\mathrm{col}(j), \mathbf{z}_c)
\]
where $\mathrm{row}(j) = j - t_{\mathrm{col}(j)}$ and
$\mathrm{col}(j) = \min_k\{t_k > j\}$.  Recovering per-chip
evaluations from $\widetilde{\mathbf{q}}$ requires a sumcheck
argument that we describe in \S\ref{sec:design-jagged-sumcheck}.

\paragraph{Implementation.}
Layout is in \texttt{crates/stark/src/jagged.rs::pack\_traces\_jagged};
the resulting \texttt{JaggedPacking\{dense\_values, chip\_infos,
offsets, total\_values, log\_dense\_size\}} carries everything the
verifier needs to reconstruct $(\mathrm{row}(j), \mathrm{col}(j))$.
The chip-info hash (\texttt{hash\_chip\_infos}) is observed into the
challenger before commit to bind chip dimensions to the transcript.

\subsection{Polynomial Commitment: WHIR}
\label{sec:design-whir}

WHIR~\cite{whir2024} is a multilinear PCS based on Reed--Solomon
proximity testing.  A WHIR commitment to $\mathbf{q} \in
\mathbb{F}^{2^n}$ encodes $\mathbf{q}$'s coefficients into an RS
codeword on a coset of size $2^{n+r}$ ($r$ = log-inv-rate, set to $4$
in our deployment for $80$-bit security) and merkle-commits the row
hashes of the codeword matrix.

The WHIR open phase combines:
\begin{enumerate}
  \item Sumcheck folding rounds (one per WHIR round, each round
        folds $f$ variables for folding-factor $f$).
  \item STIR queries at $K$ random codeword positions, with
        per-round merkle paths.
  \item Out-of-domain (OOD) sample evaluations to bind the
        polynomial against the codeword.
  \item Per-round proof-of-work (PoW) grinding.
  \item A final-round in-the-clear polynomial of degree
        $2^{\mathrm{final\_sumcheck\_rounds}}$.
\end{enumerate}

\paragraph{Folding-factor scaling.}  Our jagged shards have
$2^n \in \{2^{22}, \ldots, 2^{27}\}$, but KoalaBear's two-adicity
caps at $24$.  We scale the folding factor as
$f = \max(4, n + r - 24)$ to keep the folded codeword in the
two-adic subgroup.  For $n = 27$: $f = 7$, giving
$\lceil n / f \rceil = 4$ rounds.

\paragraph{Implementation.}
WHIR PCS configuration is in
\texttt{crates/stark/src/whir\_config.rs}; the late-binding
adapter is in \texttt{crates/stark/src/whir\_late\_binding.rs}
(see \S\ref{sec:design-late-binding}).

\subsection{Late-Binding WHIR}
\label{sec:design-late-binding}

\paragraph{Motivation.}
The upstream Plonky3 \texttt{MultilinearPcs::commit} adapter takes
opening points at commit time, absorbing them into Fiat-Shamir
\emph{before} the column-batching challenge $\alpha$.  Our protocol
derives the row coordinates $\mathbf{r}_{\mathrm{row}}$ via GKR sumcheck
\emph{after} the commit, which the upstream adapter cannot accommodate.

\begin{construction}[Late-binding WHIR]
\label{cons:late-binding}
Phase order:
\begin{enumerate}
  \item \emph{Domain separator}: observe protocol parameters into the
        challenger.
  \item \emph{Commit phase}: build the RS codeword merkle and
        observe its root.  Sample OOD points; observe their
        evaluations into the challenger.  At this point the
        EqStatement contains only OOD constraints.
  \item \emph{(Caller-driven) post-commit point derivation}: GKR/
        zerocheck runs against the now-advanced challenger,
        producing $\mathbf{r}_{\mathrm{row}}$.
  \item \emph{Late-binding registration}: for each derived point
        $\mathbf{p}$ and value $v$, the prover invokes
        $\mathtt{add\_late\_claim}(\mathbf{p}, v)$, which evaluates
        $\widetilde{\mathbf{q}}(\mathbf{p}) = v$ (recorded into the
        EqStatement) and observes $v$ into the challenger.  The
        point is \emph{not} re-observed because it was derived from
        the post-commit challenger.
  \item \emph{Open phase}: run the standard WHIR open over the
        accumulated EqStatement (OOD claims + late claims).
\end{enumerate}
\end{construction}

\paragraph{Soundness.}  Identical to upstream WHIR.  The
post-commit ordering does not weaken the proximity guarantee because
(a) the merkle commit binds the prover to a specific polynomial
\emph{before} late points are sampled, and (b) the sumcheck and
STIR-query rounds operate over the same EqStatement that contains
both OOD and late claims, with constraint-batching $\gamma$ sampled
\emph{after} all claims are observed.

\paragraph{Multi-column variant.}  For $w$ columns, the per-column
adapter is invoked $w$ times with a shared challenger, with the
prover doing all-observes-then-all-opens (and the verifier mirroring).
This works but commits $w$ independent merkle trees, which is
wasteful at scale.  The Jagged + late-binding combined construction
in \S\ref{sec:design-jagged-late-binding} reduces this to one merkle
per shard.

\paragraph{Implementation.}
\texttt{crates/stark/src/whir\_late\_binding.rs}, built directly on
\texttt{p3-whir} internals (\texttt{CommitmentWriter},
\texttt{Verifier}, \texttt{EqStatement}); no upstream patch
required.  Tests: $24$ unit tests covering single-column, multi-
column, tampering rejection, and challenger-replay.

\subsection{Jagged + Late-Binding Reduction}
\label{sec:design-jagged-late-binding}
\label{sec:design-jagged-sumcheck}

The Jagged commitment (\S\ref{sec:design-jagged}) gives one merkle
per shard but commits $\widetilde{\mathbf{q}}$, not the per-chip
column MLEs that LogUp-GKR's leaf claims need.  The reduction:

\begin{construction}[Jagged sumcheck reduction]
\label{cons:jagged-sumcheck}
Inputs:
\begin{itemize}
  \item Dense polynomial $\widetilde{\mathbf{q}}$ committed via
        Construction~\ref{cons:late-binding}.
  \item Per-chip row coordinates $\mathbf{r}_{\mathrm{row},c}$ derived
        from GKR.
  \item Per-chip per-column row-MLE values
        $y_{c,j} = \widetilde{\mathbf{T}_c[:,j]}(\mathbf{r}_{\mathrm{row},c})$
        provided by the prover.
\end{itemize}

Steps:
\begin{enumerate}
  \item Prover observes all $y_{c,j}$ into the challenger.
  \item Verifier samples batching challenge $\gamma$.
  \item Prover and verifier compute initial claim
        $T = \sum_{c,j} \gamma^{\mathrm{idx}(c,j)} \cdot y_{c,j}$.
  \item Verifier defines the weight polynomial
        $\widetilde{w}(j) = \gamma^{\mathrm{idx}(c,j)} \cdot
        \mathsf{eq}(\mathrm{row}(j), \mathbf{r}_{\mathrm{row},c})$
        for $j$ in chip $c$ column-$j$'s range, $0$ elsewhere.
  \item Run sumcheck on $T = \sum_j \widetilde{\mathbf{q}}(j) \cdot
        \widetilde{w}(j)$ over $n$ variables (degree-$2$ rounds).
        After sumcheck, the verifier has
        $\widetilde{\mathbf{q}}(\mathbf{z}^*) \cdot
         \widetilde{w}(\mathbf{z}^*) = T_{\mathrm{final}}$
        where $\mathbf{z}^*$ is the sumcheck-derived point.
  \item Prover supplies $v^* = \widetilde{\mathbf{q}}(\mathbf{z}^*)$;
        verifier independently computes $\widetilde{w}(\mathbf{z}^*)$
        from the packing metadata, checks
        $v^* \cdot \widetilde{w}(\mathbf{z}^*) = T_{\mathrm{final}}$.
  \item $v^*$ is bound to the WHIR commit by registering
        $(\mathbf{z}^*, v^*)$ as a late claim
        (Construction~\ref{cons:late-binding}, step~4).
\end{enumerate}
\end{construction}

\begin{remark}[Convention reversal]
The sumcheck protocol's eval point is in last-variable-first order
$[r_n, r_{n-1}, \ldots, r_1]$ (because pair-consecutive folding
binds the last variable first), whereas
\texttt{p3\_multilinear\_util::Poly::eval\_base} expects MSB-first
$[x_0, x_1, \ldots, x_{n-1}]$.  The $\mathbf{z}^*$ passed to
\texttt{add\_late\_claim} is reversed; \texttt{verify\_jagged\_reduction}
mirrors this on the verifier side.  This was a real integration
bug found and fixed during implementation.
\end{remark}

\paragraph{Cost analysis.}  For a fibonacci shard ($N = 19$ chips,
$\sim 30$ avg cols, $n = 27$ packed):
\begin{itemize}
  \item Commit: $1$ merkle on $2^{27+r}$ codeword vs.
        $570$ merkles on $2^{17+r}$.
  \item Open: $1$ WHIR proof on $n=27$ vs. $570$ WHIR proofs on $n=17$.
  \item Sumcheck reduction: $n = 27$ rounds, degree-$2$
        polynomials, parallelisable per round.
  \item Per-shard proof bytes: $\sim 488$\,KB vs.~$\sim 29.5$\,MB
        ($60\times$ smaller).
\end{itemize}

Per-phase profile (parallelised, fibonacci shard):
\begin{tabular}{lr}
\toprule
WHIR commit on packed dense  & $1.7$\,s ($1.9\%$) \\
Sumcheck reduction (rayon)   & $0.16$\,s ($0.2\%$) \\
WHIR open phase              & $89.4$\,s ($97.9\%$) \\
\bottomrule
\end{tabular}

\paragraph{Implementation.}
\texttt{crates/stark/src/jagged\_late\_binding.rs}, $\sim 1200$ lines
including $11$ unit tests.

\subsection{Cross-Binding}
\label{sec:design-cross-binding}

\paragraph{Threat model.}  In the dual-commit configuration (FRI at
$\zeta$ + WHIR via Construction~\ref{cons:late-binding}), a
malicious prover could in principle commit to two different traces
$\mathbf{T}_A \neq \mathbf{T}_B$ --- the FRI side via the existing
\texttt{pcs.commit} on $\mathbf{T}_A$, the WHIR side via Jagged
late-binding on $\mathbf{T}_B$.  Both commits would individually
verify, but the proof would fail to attest to a single coherent
trace.

\begin{construction}[Brakedown-style cross-binding spot check]
\label{cons:cross-binding}
After both commits are observed:
\begin{enumerate}
  \item Sample $K$ random trace positions $i_1, \ldots, i_K$ from
        the post-both-commits challenger ($K = 80$ for $\sim 80$-bit
        soundness against single-position forgery).
  \item For each $i_j$, the prover provides:
        \begin{enumerate}
          \item The trace value $\mathbf{T}[i_j]$.
          \item Merkle paths against the FRI MMCS at the codeword
                position(s) corresponding to $i_j$.
          \item Merkle paths against the WHIR MMCS at the codeword
                position corresponding to $i_j$.
        \end{enumerate}
  \item Verifier checks both paths via \texttt{Mmcs::verify\_batch}
        and that both reveal the same value.
\end{enumerate}
\end{construction}

\paragraph{Soundness.}  If $\mathbf{T}_A \neq \mathbf{T}_B$, they
differ at $\geq 1$ position; $K$ random positions catch the
divergence with probability $1 - (1 - 1/N)^K \approx 1 - 2^{-K}$
for $N = 2^{17}, K = 80$.

\paragraph{Status.}  Both sides fully implemented at host level:
prover populates FRI- and WHIR-side merkle paths
(\texttt{populate\_fri\_paths},
\texttt{prove\_cross\_binding\_spot\_check}); verifier checks both
via \texttt{Mmcs::verify\_batch}
(\texttt{verify\_cross\_binding\_spot\_check\_fri},
\texttt{verify\_cross\_binding\_spot\_check\_whir}).  The
verifier dispatcher
\texttt{verify\_cross\_binding\_for\_kb} extracts the WHIR merkle
root from the jagged bundle's \texttt{WhirProof.initial\_commitment}
and reconstructs the codeword shape via the public helper
\texttt{whir\_codeword\_dims(log\_dense\_size)}, which mirrors
\texttt{make\_chip\_late\_binding\_pcs}'s folding-factor derivation.
Tamper rejection at the merkle layer is validated by
\texttt{cross\_binding\_whir\_rejects\_tampered\_path}.

\paragraph{Implementation.}
Helpers in
\texttt{crates/stark/src/jagged\_late\_binding.rs}
(\texttt{populate\_fri\_paths},
\texttt{verify\_cross\_binding\_spot\_check\_fri},
\texttt{verify\_cross\_binding\_spot\_check\_whir},
\texttt{prove\_cross\_binding\_spot\_check}); the codeword-dimension
helper \texttt{whir\_codeword\_dims} in
\texttt{crates/stark/src/whir\_late\_binding.rs}; prover wiring in
\texttt{crates/stark/src/prover.rs}
(\texttt{try\_compute\_cross\_binding\_fri},
\texttt{prove\_jagged\_late\_binding\_with\_cross\_binding});
verifier wiring in \texttt{crates/stark/src/verifier.rs}
(\texttt{try\_verify\_cross\_binding},
\texttt{verify\_cross\_binding\_for\_kb}).

\subsection{Soundness Layers Summary}

The combined pipeline composes the following soundness layers:

\begin{tabular}{@{}lll@{}}
\toprule
\textbf{Layer}                       & \textbf{Mechanism}            & \textbf{Soundness Bound} \\
\midrule
Transition constraints               & Zerocheck (sumcheck)          & $d \cdot m / |\mathbb{F}_q|$ \\
Lookup arguments                     & LogUp-GKR + per-layer sumcheck & $3 \cdot \log_2 N / |\mathbb{F}_q|$ \\
Row-MLE binding to commit            & Late-binding WHIR             & WHIR proximity \\
FRI $\leftrightarrow$ WHIR coherence  & Cross-binding spot-check (C.1) & $1 - (1 - 1/N)^K$ \\
\bottomrule
\end{tabular}

For our parameters ($\mathbb{F}_q \approx 2^{124}$, $m \leq 27$,
$d = 2$, $N \leq 2^{27}$, $K = 80$), each layer's soundness error is
$\leq 2^{-80}$.

\subsection{Implementation Activation Matrix}

\begin{tabular}{@{}ll@{}}
\toprule
\textbf{Configuration}                                  & \textbf{Env flags} \\
\midrule
FRI (default, unchanged)                                & \emph{(none)} \\
WHIR fast path + per-chip late-binding                  & \texttt{ZIREN\_USE\_WHIR=1} \\
WHIR fast path + jagged late-binding                    & \texttt{ZIREN\_USE\_WHIR=1
                                                            ZIREN\_LATE\_BINDING=jagged} \\
+ cross-binding (FRI + WHIR sides)                      & + \texttt{ZIREN\_CROSS\_BINDING=1} \\
\bottomrule
\end{tabular}

Production deployment should use the fourth row (jagged late-binding
with cross-binding).  End-to-end smoke tests on fibonacci, JSON,
keccak, and large-sum (multi-shard) workloads validate the first
three rows; the fourth row's merkle-layer correctness is validated
by lib-level round-trip and tamper-rejection unit tests against
the same MMCS instance the prover uses.


% =============================================================================
% Section 7: Phase 2a Construction --- Zerocheck for Transition Constraints
\section{Phase 2a: Zerocheck Integration}
\label{sec:zerocheck-phase2a}

Phase 2a delivers the sumcheck infrastructure that replaces the quotient
polynomial in the WHIR fast path.  The integration is complete and
end-to-end verified for fibonacci; the remaining soundness gap (transition
constraints not yet bound to a PCS opening) is scheduled for the follow-up
phase.

\subsection{What Ships}

\paragraph{Sumcheck core (\texttt{crates/stark/src/zerocheck\_prover.rs}, 5 unit tests, all passing).}
\begin{itemize}
  \item \texttt{eq\_mle\_table(r)} -- dense evaluation of $\mathsf{eq}(r,\cdot)$ over $\{0,1\}^m$.
  \item \texttt{MultilinearExt::\{fold\_first, evaluate\}} -- standard MLE operations.
  \item \texttt{prove\_zerocheck\_multilinear(c\_evals, r)} -- deterministic prover.
  \item \texttt{verify\_zerocheck\_rounds(proof, r, initial\_claim)} -- deterministic verifier.
  \item \texttt{prove\_zerocheck\_with\_challenger(c\_evals, m, challenger)} -- Fiat-Shamir prover.
  \item \texttt{verify\_zerocheck\_with\_challenger(proof, m, initial\_claim, challenger)} -- Fiat-Shamir verifier.
\end{itemize}

\paragraph{Hypercube constraint evaluator.}
\texttt{eval\_constraints\_on\_hypercube} iterates each row of the Boolean
hypercube, lifts the (base-field) main and preprocessed rows to the
extension field, and invokes the chip's \texttt{Air::eval} via
\texttt{VerifierConstraintFolder}.  Produces a length-$2^m$ table of
extension-field values representing
$$C_{\mathrm{batched}}(b) = \sum_j \alpha^j \cdot C_j\bigl(\mathrm{row}_b, \mathrm{row}_{b+1}, \mathrm{preproc}_b, \mathrm{preproc}_{b+1}, \mathrm{pub}\bigr).$$

Selectors are set analytically per row ($\texttt{is\_first\_row}$,
$\texttt{is\_last\_row}$, $\texttt{is\_transition}$); next-row accesses use
cyclic wrap within the hypercube.

\paragraph{\texttt{ShardProof::zerocheck\_proofs} field.}
New optional field serialized with \texttt{\#[serde(default)]} for
backward-compatibility:
\begin{verbatim}
pub struct ShardProof<SC: StarkGenericConfig> {
    // existing fields
    #[serde(default)]
    pub zerocheck_proofs: Option<Vec<ZerocheckProof<Challenge<SC>>>>,
}
\end{verbatim}

\paragraph{Prover wiring (WHIR fast path).}
After sampling $\alpha$:
\begin{itemize}
  \item If \texttt{chip.permutation\_width() > 0} (lookup-dependent chips):
        emit an empty placeholder proof.  These chips become eligible for
        zerocheck once Phase~2b removes the permutation argument.
  \item Otherwise: run \texttt{eval\_constraints\_on\_hypercube} and
        \texttt{prove\_zerocheck\_with\_challenger}.
\end{itemize}
Timing is logged to \texttt{/tmp/ziren\_open\_breakdown.txt} under
\texttt{ZEROCHECK total=...ms}.

\paragraph{Verifier wiring.}
After sampling $\alpha$, if \texttt{zerocheck\_proofs} is present, for each
chip/proof pair the verifier replays the sumcheck transcript.  Empty
\texttt{rounds} means "skipped" (permutation-dependent chip).  Non-empty
proofs are passed to \texttt{verify\_zerocheck\_with\_challenger}.  New
error variants: \texttt{ZerocheckFailed}, \texttt{InvalidProofShape}.

\paragraph{Downstream fixes.}
\texttt{ShardProof} literals updated in
\texttt{crates/sdk/src/provers/mock.rs} and
\texttt{crates/recursion/circuit/src/stark.rs}.

\subsection{End-to-End Test Results}

\begin{table}[h]
\centering
\begin{tabular}{lll}
\toprule
\textbf{Mode} & \textbf{Result} & \textbf{Notes} \\
\midrule
FRI (quotient + permutation)                            & PASS (8.1\,s) & Existing pipeline unchanged \\
WHIR core (zerocheck emitted, skipped for all 19 chips) & PASS (6.8\,s) & Serialization round-trip works \\
WHIR compressed                                          & PASS (6.9\,s) & Recursion verifier bypasses empty proofs \\
\bottomrule
\end{tabular}
\caption{Phase 2a end-to-end verification on fibonacci.}
\label{tab:phase2a-e2e}
\end{table}

On fibonacci, all 19 CPU-program chips have \texttt{permutation\_width() > 0},
so \texttt{ZEROCHECK total=0ms chips=19} --- the sumcheck is never actually
executed for this workload.  Real-world exercise of the zerocheck math
will come once Phase~2b removes the permutation argument.

\subsection{Phase 2a Scope Limitation (By Design)}

Chips that use the lookup permutation argument are skipped because their
\texttt{Air::eval} calls \texttt{builder.permutation()}, which requires a
permutation trace we don't have in the WHIR fast path.  Two ways to lift
this:

\begin{enumerate}
  \item \textbf{Phase 2b} (LogUp-GKR): replace the permutation argument
        with GKR.  Chips no longer need \texttt{builder.permutation()}.
        All chips become zerocheck-eligible.
  \item \textbf{Interim fallback}: generate the permutation trace in the
        WHIR fast path purely for the zerocheck eval (not committed).
\end{enumerate}

Phase~2b is the correct long-term fix and matches SP1's hypercube
architecture.

\subsection{Files Changed (Phase 2a)}

\begin{verbatim}
crates/stark/src/lib.rs                    +1  module declaration
crates/stark/src/zerocheck.rs              +1  derive Serialize/Deserialize
crates/stark/src/zerocheck_prover.rs       +420 NEW
crates/stark/src/types.rs                  +5   ShardProof field
crates/stark/src/prover.rs                 +55  WHIR path wiring
crates/stark/src/verifier.rs               +40  zerocheck verify + errors
crates/sdk/src/provers/mock.rs             +1
crates/recursion/circuit/src/stark.rs      +1
\end{verbatim}

\section{Phase 2b: LogUp-GKR Protocol Core}
\label{sec:logup-gkr-phase2b}

Phase 2b Step 1 delivers a complete, unit-tested LogUp-GKR fraction-sum
protocol in \texttt{crates/stark/src/logup\_gkr.rs}.  The protocol takes a
set of leaf fractions $\{(\text{num}_i, \text{denom}_i)\}$ and lets the
prover commit to the claim
$$\sum_i \frac{\text{num}_i}{\text{denom}_i} = 0$$
via a logarithmic-depth reduction.  Each GKR layer is a \emph{folding}
step (not a sumcheck round) because the fraction-combine identity is
bilinear in the two halves.

\subsection{Fraction-Sum Combine}

For two fractions $(a, b) = a/b$ and $(c, d) = c/d$,
$$
\frac{a}{b} + \frac{c}{d} = \frac{a \cdot d + b \cdot c}{b \cdot d}.
$$
We encode this as the binary operation
$$(a,b) \oplus (c,d) = (a \cdot d + b \cdot c,\; b \cdot d).$$

The prover builds a balanced binary tree over leaf fractions using
$\oplus$; the root is $(N_m, D_m)$.  For a balanced LogUp argument,
$N_m = 0$, which the verifier checks in the clear.

\subsection{GKR Reduction}

At each layer $k$ (from root $m$ down to leaves $0$), the current claim
is a multilinear evaluation of $(N_k, D_k)$ at a point $z \in \mathbb{F}_{\mathrm{ext}}^{m-k}$.
The identity
\begin{align*}
N_{k-1}(z \Vert 0) \cdot D_{k-1}(z \Vert 1) + D_{k-1}(z \Vert 0) \cdot N_{k-1}(z \Vert 1) &= N_k(z), \\
D_{k-1}(z \Vert 0) \cdot D_{k-1}(z \Vert 1) &= D_k(z),
\end{align*}
lets the prover produce four values $(N_k(z\Vert 0), N_k(z\Vert 1), D_k(z\Vert 0), D_k(z\Vert 1))$
that the verifier checks against the current claim via the equations
above.  The verifier samples $r^* \in \mathbb{F}_{\mathrm{ext}}$ and
folds the claim to the new point $z \Vert r^*$:
\begin{align*}
N_{k-1}(z \Vert r^*) &= (1 - r^*) \cdot N_{k-1}(z \Vert 0) + r^* \cdot N_{k-1}(z \Vert 1), \\
D_{k-1}(z \Vert r^*) &= (1 - r^*) \cdot D_{k-1}(z \Vert 0) + r^* \cdot D_{k-1}(z \Vert 1).
\end{align*}

After $m$ layers, the remaining claim is
$\bigl(N_0(\mathsf{eval\_point}), D_0(\mathsf{eval\_point})\bigr)$ -- the
multilinear evaluation of the \emph{leaf} fractions at the accumulated
point.  The caller closes the loop by checking this against the
fingerprints $(\text{multiplicity}_i,\; \alpha - \text{fingerprint}_i)$
reconstructed from the main-trace PCS opening at \texttt{eval\_point}.

\subsection{Soundness Tests}

Six unit tests validate the protocol:

\begin{table}[h]
\centering
\begin{tabular}{ll}
\toprule
\textbf{Test} & \textbf{What it establishes} \\
\midrule
\texttt{fraction\_combine\_is\_associative\_on\_small\_input} & Algebraic sanity of $\oplus$ \\
\texttt{sum\_of\_balanced\_fractions\_has\_zero\_numerator}   & Target identity: balanced sends/receives $\Rightarrow N_m = 0$ \\
\texttt{logup\_gkr\_prove\_verify\_roundtrip\_on\_balanced\_input} & Honest prover round-trip \\
\texttt{logup\_gkr\_leaf\_claim\_matches\_mle\_evaluation}    & Reduction identity: final claim = $N_0(\mathsf{eval\_point})$ \\
\texttt{logup\_gkr\_rejects\_tampered\_combine\_identity}     & Soundness: corrupted combine round is rejected \\
\texttt{logup\_gkr\_rejects\_tampered\_root}                  & Soundness: corrupted root is rejected \\
\bottomrule
\end{tabular}
\caption{LogUp-GKR unit tests, all passing.}
\label{tab:logup-gkr-tests}
\end{table}

The \texttt{leaf\_claim\_matches\_mle\_evaluation} test is the key
soundness property: without it, the protocol might appear to verify but
would not actually bind the prover to the committed fraction leaves.

\subsection{Public API}

\begin{verbatim}
pub struct Fraction<EF>      { pub num: EF, pub denom: EF }
pub struct LogUpGkrProof<EF> {
    pub root: (EF, EF),
    pub rounds: Vec<(EF, EF, EF, EF)>,  // (n0, n1, d0, d1) per layer
    pub eval_point: Vec<EF>,
    pub leaf_claim: (EF, EF),
}

pub fn build_fraction_tree<EF>(leaves: &[Fraction<EF>]) -> Vec<Vec<Fraction<EF>>>;
pub fn sum_of_fractions<EF>(leaves: &[Fraction<EF>]) -> Fraction<EF>;
pub fn eval_mle<EF>(table: &[EF], r: &[EF]) -> EF;
pub fn eval_mle_first_var<EF>(table: &[EF], r: &[EF]) -> EF;

pub fn prove_logup_gkr<F, EF, Challenger>(
    leaves: &[Fraction<EF>], challenger: &mut Challenger
) -> LogUpGkrProof<EF>;

pub fn verify_logup_gkr<F, EF, Challenger>(
    proof: &LogUpGkrProof<EF>, challenger: &mut Challenger
) -> Option<(Vec<EF>, Fraction<EF>)>;
\end{verbatim}

The \texttt{eval\_point} returned is in MSB-first order (matching
\texttt{eval\_mle\_first\_var}), so the caller can evaluate the leaf MLE
directly at that point with no re-ordering.

\subsection{Integration Path (Phase 2b Steps 2 -- 4)}

Outstanding work to connect this protocol to the native prover:

\begin{enumerate}
  \item \textbf{Build the leaves from lookup events}. For each chip:
        \begin{itemize}
          \item Enumerate send/receive events;
          \item For each event, compute the fingerprint
                $f = \beta \cdot \text{payload}_0 + \beta^2 \cdot \text{payload}_1 + \dots$;
          \item Leaf numerator $= \pm \text{multiplicity}$
                ($+$ for sends, $-$ for receives);
          \item Leaf denominator $= \alpha - f$.
        \end{itemize}
  \item \textbf{Attach the GKR proof to \texttt{ShardProof}}.  Add a new
        optional field
        \texttt{logup\_gkr\_proofs: Option<Vec<LogUpGkrProof<Challenge<SC>>>>}
        with \texttt{\#[serde(default)]}.
  \item \textbf{Remove the permutation trace in WHIR mode}.  Skip
        \texttt{generate\_permutation\_trace} when LogUp-GKR is enabled.
        Drop the permutation-commit round from the Fiat-Shamir transcript.
  \item \textbf{Verifier: close the loop}.  After verifying the GKR layers,
        check the leaf claim against fingerprints recomputed from the
        opened trace values at \texttt{eval\_point}.
  \item \textbf{Unblock zerocheck}.  Once the permutation argument is
        gone, every chip becomes zerocheck-eligible; the Phase~2a
        limitation at \texttt{chip.permutation\_width() > 0} can be
        dropped.
\end{enumerate}

\subsection{Combined Soundness Picture}

\begin{table}[h]
\centering
\small
\begin{tabular}{lcccc}
\toprule
\textbf{Component} & \textbf{FRI} & \textbf{WHIR (pre-2a)} & \textbf{WHIR (2a only)} & \textbf{WHIR (2a + 2b + multi-point open)} \\
\midrule
Main trace binding    & \checkmark & \checkmark & \checkmark & \checkmark \\
Lookup integrity      & \checkmark & $\bigtriangleup$\footnotemark[1] & $\bigtriangleup$\footnotemark[1] & \checkmark (GKR) \\
Transition constraints & \checkmark & $\times$ & $\bigtriangleup$\footnotemark[2] & \checkmark (sumcheck) \\
Production-sound?     & \checkmark & $\times$ & $\times$ & \checkmark \\
\bottomrule
\end{tabular}
\caption{Soundness status across phases.}
\label{tab:soundness-phases}
\end{table}
\footnotetext[1]{Only the cumulative-sum balance is checked.}
\footnotetext[2]{Sumcheck transcript is bound, but the final claim is not yet checked against a PCS opening.}


% =============================================================================
% Section 7: Phase 2b Construction --- LogUp-GKR with Sumcheck Reduction
\section{LogUp-GKR: Bug, Fix, and End-to-End Integration}
\label{sec:logup-gkr-fix}

\subsection{The Soundness Bug (And Why It Was Hard to Spot)}

The initial Phase 2b implementation in
\texttt{crates/stark/src/logup\_gkr.rs} used a folding-only reduction.
At each GKR layer, the prover sent four values
\[
  (N_k(z \Vert 0),\; N_k(z \Vert 1),\; D_k(z \Vert 0),\; D_k(z \Vert 1)),
\]
and the verifier checked the combine identity
\begin{align*}
  N_{k+1}(z) &\stackrel{?}{=} N_k(z \Vert 0) \cdot D_k(z \Vert 1) + D_k(z \Vert 0) \cdot N_k(z \Vert 1), \\
  D_{k+1}(z) &\stackrel{?}{=} D_k(z \Vert 0) \cdot D_k(z \Vert 1),
\end{align*}
before folding a challenge $r^*$ to obtain the next claim at $z \Vert r^*$.

This identity is only valid for Boolean $z$. At a random $z \in \mathbb{F}^{m-k-1}$
the right-hand side is a degree-2 polynomial in each coordinate, while
$N_{k+1}$ is multilinear (degree 1). For arbitrary $z$ the identity fails.

\paragraph{Why the original test suite didn't catch it.}
The honest-roundtrip test used balanced fractions with all denominators
drawn from $\{7, 11\}$. Those repeated denominators made the degree-2
cross-terms vanish algebraically, so the check accidentally held. A
regression test with varied denominators
(\texttt{logup\_gkr\_roundtrip\_with\_varied\_denominators}) fails with
the folding-only reduction but now passes with the sumcheck-based fix.

\subsection{The Correct Reduction}

At each GKR layer we need to reduce a claim
$\bigl(N_{k+1}(z),\; D_{k+1}(z)\bigr)$ into claims on $N_k, D_k$ at a
fresh point. The MLE-extended combine identity
\[
  \widetilde{N}_{k+1}(z)
   \;=\; \sum_{b \in \{0,1\}^{m-k-1}} \mathsf{eq}(z, b) \cdot
         \bigl[\widetilde{N}_k(b, 0) \cdot \widetilde{D}_k(b, 1)
               + \widetilde{D}_k(b, 0) \cdot \widetilde{N}_k(b, 1)\bigr]
\]
is genuinely identical at every $z$ because the product of multilinear
functions is only evaluated at \emph{Boolean} $b$ inside the sum. We
turn this into a sumcheck claim.

\subsubsection*{Protocol per layer}
\begin{enumerate}
  \item \textbf{Batch.} Sample $\lambda \in \mathbb{F}_{\mathrm{ext}}$. The
        merged claim is $\lambda \cdot \widetilde{N}_{k+1}(z) + \widetilde{D}_{k+1}(z)$.
  \item \textbf{Sumcheck.} Run $m-k-1$ rounds of sumcheck over
        \[
          g(b) \;=\; \mathsf{eq}(z, b) \cdot
               \bigl(\lambda \cdot (\widetilde{N}_k \widetilde{D}_k \text{ cross})
                     + \widetilde{D}_k(b, 0) \cdot \widetilde{D}_k(b, 1)\bigr).
        \]
        $g$ has total degree 3 in each variable being summed out, so each
        round sends 4 field elements (evaluations at $X \in \{0, 1, 2, 3\}$).
  \item \textbf{Final evaluation.} The prover sends
        $(\widetilde{N}_k(r^*, 0),\; \widetilde{N}_k(r^*, 1),\;
          \widetilde{D}_k(r^*, 0),\; \widetilde{D}_k(r^*, 1))$. The
        verifier checks the sumcheck's final identity directly.
  \item \textbf{Line protocol.} Sample $t \in \mathbb{F}_{\mathrm{ext}}$
        and fold the two pairs of values into a single claim at $(r^*, t)$.
\end{enumerate}

After $m$ layers, the accumulated point $\mathbf{r} \in \mathbb{F}^m$ and
the final claim $(\widetilde{N}_0(\mathbf{r}), \widetilde{D}_0(\mathbf{r}))$
can be checked against the fingerprints reconstructed from the
main-trace PCS opening at $\mathbf{r}$.

\subsection{Implementation}

\paragraph{Proof shape.} The new \texttt{LogUpGkrProof} carries per-layer
reductions:
\begin{verbatim}
pub struct LogUpGkrLayerProof<EF> {
    pub sumcheck_rounds: Vec<[EF; 4]>, // degree-3 round polys
    pub final_evals: [EF; 4],          // (N_k(r*,0), N_k(r*,1), D_k(r*,0), D_k(r*,1))
}

pub struct LogUpGkrProof<EF> {
    pub root: (EF, EF),
    pub layers: Vec<LogUpGkrLayerProof<EF>>,
    pub eval_point: Vec<EF>,
    pub leaf_claim: (EF, EF),
}
\end{verbatim}

\paragraph{Variable ordering.} The sumcheck folds the last variable
first at every round (cache-friendly pair-consecutive layout), so the
raw \texttt{r\_star} vector built during the sumcheck encodes challenges
in reverse variable order. The prover and verifier both reverse
\texttt{r\_star} before appending the line challenge, so the returned
\texttt{eval\_point} matches the \texttt{eval\_mle\_first\_var}
convention used by the caller.

\paragraph{Unit tests.} All 8 tests pass:
\begin{itemize}
  \item \texttt{fraction\_combine\_is\_associative\_on\_small\_input}
  \item \texttt{sum\_of\_balanced\_fractions\_has\_zero\_numerator}
  \item \texttt{logup\_gkr\_prove\_verify\_roundtrip\_on\_balanced\_input}
  \item \texttt{logup\_gkr\_leaf\_claim\_matches\_mle\_evaluation}
        \textit{(final claim equals MLE at \texttt{eval\_point})}
  \item \texttt{logup\_gkr\_roundtrip\_with\_varied\_denominators}
        \textit{(regression for the folding-only bug)}
  \item \texttt{logup\_gkr\_rejects\_tampered\_final\_evals}
  \item \texttt{logup\_gkr\_rejects\_tampered\_root}
  \item \texttt{logup\_gkr\_leaf\_claim\_reconstructs\_from\_row\_mle}
        \textit{(end-to-end: real chip $\to$ leaves $\to$ GKR $\to$
        reconstruction from row-MLE opening matches)}
\end{itemize}

\paragraph{End-to-end.} The WHIR fast path now emits LogUp-GKR proofs
per chip and the verifier replays them. Fibonacci core proof: 7.2\,s
(1.15\,kHz); compressed: 7.1\,s.

\begin{table}[h]
\centering
\begin{tabular}{lrrl}
\toprule
\textbf{Mode} & \textbf{e2e (s)} & \textbf{kHz} & \textbf{Status} \\
\midrule
main FRI            & 12.54 & 0.66 & Baseline \\
upgrade FRI         & 8.21  & 1.01 & +53\% vs main \\
upgrade WHIR (zerocheck skipped, no GKR) & 6.8  & 1.22 & +85\% vs main, unsound \\
upgrade WHIR + sumcheck LogUp-GKR        & 7.2  & 1.15 & +74\% vs main, transcript-bound GKR \\
\bottomrule
\end{tabular}
\caption{Fibonacci core-proof times across configurations. LogUp-GKR
adds ~5\% over the bare WHIR path thanks to the 1.5\,s per-shard sumcheck
cost. Compared to main FRI, we keep 74\% of the measured speedup while
restoring lookup-argument binding.}
\label{tab:logup-gkr-perf}
\end{table}

\subsection{Closing the Soundness Loop}

The \texttt{reconstruct\_leaf\_claim\_from\_openings} helper
(\texttt{crates/stark/src/logup\_gkr.rs}) implements the final soundness
check.  Given the opened main-trace and preprocessed-trace values at the
row coordinates of $\mathbf{r}$, it evaluates each lookup interaction's
multiplicity and fingerprint expressions at that row point (using
\texttt{VirtualPairCol::apply::<EF, EF>}), then combines the
per-interaction fractions into the leaf-layer MLE values at
$\mathbf{r}$:
\begin{align*}
  \widetilde{N}_0(\mathbf{r})
    &= \sum_{j} \mathsf{eq}(\mathbf{r}_{\mathrm{int}}, j) \cdot \mathrm{mult}_j(\mathbf{r}_{\mathrm{row}}), \\
  \widetilde{D}_0(\mathbf{r})
    &= \sum_{j} \mathsf{eq}(\mathbf{r}_{\mathrm{int}}, j) \cdot \bigl(\alpha + \beta \cdot \mathrm{kind}_j + \textstyle\sum_i \beta^{i+1} \cdot \mathrm{value}_{j,i}(\mathbf{r}_{\mathrm{row}})\bigr).
\end{align*}

The unit test \texttt{logup\_gkr\_leaf\_claim\_reconstructs\_from\_row\_mle}
builds a real single-column chip with a single lookup, runs the GKR
prover, evaluates the row-MLE of the trace at \texttt{proof.eval\_point}
using \texttt{eval\_mle\_first\_var}, calls the reconstruction helper,
and confirms the result matches \texttt{proof.leaf\_claim} exactly ---
closing the protocol's soundness chain end-to-end on synthetic data.

\paragraph{Outstanding infrastructure.} What the verifier still lacks is
the ability to \emph{obtain} the main-trace row-MLE values at
\texttt{proof.eval\_point}: today the WHIR fast path opens the trace
only at a univariate point $\zeta$.  Once the PCS layer exposes
multilinear opening at arbitrary points (the same infrastructure needed
to close the Phase~2a zerocheck gap), the verifier will:
\begin{enumerate}
  \item After \texttt{verify\_logup\_gkr}, obtain the opened trace
        values at \texttt{proof.eval\_point} via WHIR.
  \item Split \texttt{eval\_point} into \texttt{(r\_row, r\_int)}.
  \item Call \texttt{reconstruct\_leaf\_claim\_from\_openings}.
  \item Compare the reconstruction to \texttt{proof.leaf\_claim}.
\end{enumerate}
All of the arithmetic is already implemented and unit-tested; only the
PCS-level plumbing remains.

\subsection{Files Touched}

\begin{verbatim}
crates/stark/src/logup_gkr.rs         rewritten with sumcheck reduction
crates/stark/src/zerocheck_prover.rs  made fold_table_first pub
crates/stark/src/prover.rs            WHIR path emits LogUpGkrProof per chip
crates/stark/src/verifier.rs          WHIR path replays LogUpGkrProof
crates/stark/src/types.rs             ShardProof::logup_gkr_proofs field
crates/sdk/src/provers/mock.rs        literal update
crates/recursion/circuit/src/stark.rs literal update
\end{verbatim}


% =============================================================================
% Section 8: Phase 2c Construction --- Late-Binding WHIR + Jagged + Cross-Binding
\section{Phase 2c: WHIR Late-Binding and Per-Chip Row-MLE}
\label{sec:whir-late-binding}

Phase 2c closes the soundness gap left by Phases 2a (zerocheck) and 2b
(LogUp-GKR): the per-chip \texttt{LogUpRowOpening} values are bound
cryptographically to the main-trace WHIR commitment.  This requires
two pieces of new infrastructure: a WHIR PCS adapter that supports
\emph{late-binding} evaluation points, and a restructured GKR leaf
claim that needs only a single $\alpha$-combined row-MLE value per
chip rather than one per column.

\subsection{The Problem with the Standard WHIR Adapter}

The upstream \texttt{MultilinearPcs::commit(poly, opening\_points, ch)}
adapter requires the evaluation points at commit time: opening points
are absorbed into the Fiat-Shamir transcript \emph{before} the
column-batching challenge $\alpha$ is sampled.  This binding order is
essential for the upstream soundness proof.

Our protocol cannot satisfy that ordering.  After committing the main
trace we run zerocheck (which derives $\zeta$) and LogUp-GKR (which
derives a per-chip evaluation point $r_{\mathrm{row}}$); only then do
we know the full set of points at which we want to open the trace.
Sampling $r_{\mathrm{row}}$ before commit is impossible because GKR's
sumcheck transcript depends on the merkle root.

\subsection{Late-Binding Adapter}

\paragraph{Insight.}
WHIR's commitment phase already adds OOD constraints to the equality
statement \emph{after} the merkle root is observed
(\texttt{committer.commit} in \texttt{p3-whir}).  The cryptography
permits adding constraints between merkle commit and the WHIR rounds;
only the upstream API hides this capability behind a one-shot
\texttt{commit} call.

\paragraph{Implementation.}
\texttt{crates/stark/src/whir\_late\_binding.rs} (single column):
\begin{itemize}
  \item \texttt{commit\_column(poly, ch) -> (commit, prover\_state)}\\
        domain separator $\to$ width $\to$ merkle commit, observe root
        $\to$ OOD sample/observe.
  \item \texttt{add\_late\_claim(state, point, ch) -> value}\\
        evaluate $p(\mathrm{point})$, add to statement, observe
        \emph{value} into challenger.  Point itself is not re-observed:
        it is already bound by the caller's downstream protocol that
        derived it.
  \item \texttt{open(state, ch) -> (proof, late\_claims)}\\
        runs WHIR's prove phase against the accumulated statement.
  \item \texttt{replay\_commit\_into\_challenger(proof, ch) -> parsed}\\
        verifier-side: re-runs the commit-phase challenger
        interactions to reach the same state the prover had.
  \item \texttt{finish\_verify(proof, parsed, late\_claims, ch) -> Result}\\
        verifier-side: builds the equality statement in the same order
        the prover did (OOD first, then late claims), passes an
        empty-OOD \texttt{ParsedCommitment} to suppress the internal
        \texttt{concatenate} (which would re-add OOD constraints in
        the wrong position and corrupt $\gamma$-batching).
\end{itemize}

\paragraph{Soundness validation.}
Two unit tests in \texttt{whir\_late\_binding.rs}:
\begin{itemize}
  \item \texttt{late\_binding\_round\_trip\_single\_column}: end-to-end
        commit $\to$ derive point $\to$ add claim $\to$ prove $\to$
        verify with prover and verifier sampling the late point from
        their own challengers (sanity-checked via \texttt{assert\_eq!}).
  \item \texttt{late\_binding\_rejects\_tampered\_value}: confirms the
        binding is cryptographic --- a tampered late-claim value is
        rejected by \texttt{finish\_verify} (returns
        \texttt{Sumcheck(InvalidPowWitness)}).
\end{itemize}

\paragraph{Constraint ordering subtlety.}
\texttt{combine\_hypercube} weights the $i$-th constraint by
$\gamma^i$.  The prover's statement is
$[\mathrm{ood}_0, \ldots, \mathrm{ood}_{N-1}, \mathrm{late}_0,
\ldots, \mathrm{late}_{L-1}]$, set by
\texttt{CommitmentWriter.commit} (OOD first) then
\texttt{add\_late\_claim} (late after).  The naive verifier would call
\texttt{Verifier::verify(parsed, public\_statement = late\_claims)},
which internally does \texttt{statement.concatenate(\&parsed.ood)} and
ends up with $[\mathrm{late}, \mathrm{ood}]$ --- different
$\gamma^i$-weights, transcript mismatch, sumcheck rejects.  Our
\texttt{finish\_verify} builds the full $[\mathrm{ood}, \mathrm{late}]$
statement explicitly and passes an empty-OOD \texttt{ParsedCommitment}
to neutralise the internal concatenation.

\subsection{Architectural Choice for Multi-Column Binding}

Each chip's main trace has $w$ columns.  We need to bind the
per-column row-MLE evaluations $v_j = \mathsf{MLE}(\mathrm{col}_j)
(r_{\mathrm{row}})$ in the chip's \texttt{LogUpRowOpening} to the
WHIR commitment of the chip's trace.

Two viable approaches:

\paragraph{(A) Per-column WHIR commits.}
Commit each of the $w$ columns separately, run the late-binding
adapter $w$ times.  Soundness clean: each $v_j$ is bound by its own
WHIR opening at $r_{\mathrm{row}}$.  Cost: $w \times \#\mathrm{chips}$
merkle trees and WHIR proofs.  For a typical RISC-V STARK
($w \approx 30$, $\#\mathrm{chips} \approx 100$) this is $\sim 3000$
WHIR commits per shard --- a $30\times$ overhead vs the existing
single-batch commit.

\paragraph{(B) Single $\alpha$-combined commit + restructured GKR.}
Commit each chip's trace as a single multi-column matrix.  WHIR
internally batches columns via $g_c = \sum_{j=0}^{w-1} \alpha_c^j
\cdot \mathrm{col}_{c,j}$ where $\alpha_c$ is sampled at commit time
for chip $c$.  Restructure LogUp-GKR so its leaf-claim reconstruction
needs only the \emph{combined} row-MLE
$V_c := g_c(r_{\mathrm{row}}) = \sum_j \alpha_c^j \cdot v_{c,j}$
rather than the individual $v_{c,j}$.  Cost: one WHIR commit per
chip, matching the existing structure.

\paragraph{Why (B) is sound.}
The malicious-prover attack against the naive multi-column late
binding is: prover sees $\alpha_c$ before $r_{\mathrm{row}}$, then
chooses per-column values $v_{c,j}$ to satisfy
$\sum_j \alpha_c^j \cdot v_{c,j} = g_c(r_{\mathrm{row}})$, gaining
$w-1$ degrees of freedom.  Approach (B) eliminates this attack by
\emph{never trusting the prover with per-column values}: only the
combined value $V_c$ is needed for verification, and $V_c$ is bound
directly by WHIR's opening proof.  No $w-1$ slack exists because
there is no $w$-dimensional unknown.

\paragraph{Why we choose (B).}
Soundness is identical to (A); cost is $30\times$ smaller; the
existing chip commit shape carries through with no per-column
explosion.  The price is a non-trivial change to the LogUp-GKR
leaf-claim reconstruction.

\subsection{Restructured Leaf-Claim Reconstruction}

Current (Phase 2b) leaf-claim shape (\texttt{logup\_gkr.rs}):
$$
(N_{\mathrm{leaf}}, D_{\mathrm{leaf}}) =
\bigoplus_i \bigl(m_i,\ \alpha - f_i(\mathrm{row})\bigr)
$$
where $m_i$ is the multiplicity and $f_i$ is the fingerprint of the
$i$-th interaction at row $\mathrm{row}$.  Each $f_i$ is a polynomial
in the per-column values $v_{c,j} = \mathsf{MLE}(\mathrm{col}_{c,j})
(r_{\mathrm{row}})$ of the chip's main and preprocessed columns
(typically $f_i(\mathrm{row}) = \beta + \mathrm{kind}_i + \sum_k
\beta^k \cdot v_{c,j_{i,k}}$).

The reconstruction
\texttt{reconstruct\_leaf\_claim\_from\_openings} currently:
\begin{enumerate}
  \item Reads $v_{c,j}$ from \texttt{LogUpRowOpening.main\_at\_r\_row}
        and \texttt{preproc\_at\_r\_row}.
  \item For each interaction $i$, computes $f_i$ from the relevant
        $v_{c,j}$.
  \item Folds the $L$ interactions through the fraction-sum tree to
        get $(N_{\mathrm{leaf}}, D_{\mathrm{leaf}})$.
\end{enumerate}

Restructured (Phase 2c) leaf-claim shape:
$$
V_c := \sum_{j=0}^{w_c-1} \alpha_c^j \cdot v_{c,j}
\quad\text{(single per-chip combined value, bound by WHIR)}
$$
$$
V_c^{\mathrm{prep}} := \sum_{j=0}^{w'_c-1} \alpha_c^{w_c+j} \cdot v_{c,j}^{\mathrm{prep}}
\quad\text{(combined preprocessed, same chip's $\alpha_c$)}
$$

The reconstruction now needs to compute $f_i$ given only $V_c$ and
$V_c^{\mathrm{prep}}$ (and the public per-interaction shape).  This
is impossible in general: $f_i$ depends on \emph{specific} columns
$v_{c,j_{i,k}}$, not their $\alpha_c$-mix.

\subsubsection{The Fix: Fingerprint Re-Batching at GKR Time}

Restructure the GKR leaves so each leaf $\ell$ is itself a single
$\beta$-combined value carrying enough information that the
verifier can reconstruct it from $V_c$ alone.

Concretely, fold the per-interaction fingerprint into the column
batching: instead of computing $f_i$ from individual $v_{c,j}$, we
choose the column-batching coefficients such that
$$
f_i = \beta_i \cdot V_c + \mathrm{(public terms)}
$$
where $\beta_i$ is fixed by the protocol's interaction shape and is
the same on both sides.  The challenge is that different interactions
$i$ touch different subsets of columns; mapping all of them into a
single $V_c$ requires redefining the column-batching coefficients
per-interaction.

\paragraph{Concrete construction.}
For chip $c$ with $L_c$ interactions and total column count
$w_c + w'_c$ (main + preprocessed), define:
\begin{itemize}
  \item Per-interaction column-projection vectors
        $\pi_i \in \mathbb{F}^{w_c + w'_c}$ where $\pi_i[j] = \beta^k$
        if interaction $i$ uses column $j$ in slot $k$, else $0$.
  \item A single chip-wide combined polynomial
        $g_c(\mathbf{x}) = \sum_i \gamma^i \cdot \langle \pi_i,
        \mathrm{cols}(\mathbf{x}) \rangle$
        where $\gamma$ is sampled from the post-merkle challenger and
        $\mathrm{cols}(\mathbf{x})$ is the column vector of the chip
        evaluated at $\mathbf{x}$.
  \item Bind $g_c(r_{\mathrm{row}})$ via WHIR late-binding.
\end{itemize}

The verifier reconstructs the leaf claim from a \emph{single}
$g_c(r_{\mathrm{row}})$ value.  The per-interaction fingerprint
contributions are stitched together via $\gamma^i$ binding, which is
sampled \emph{after} $r_{\mathrm{row}}$ is known --- so all degrees
of freedom in per-column values are eliminated by the
$\gamma$-randomness.

\paragraph{Soundness sketch.}
A malicious prover wishing to lie about a leaf-claim value would need
to produce a $g_c$ such that $g_c(r_{\mathrm{row}})$ matches a
fabricated leaf-claim.  Because $\gamma$ is sampled after
$r_{\mathrm{row}}$ and after the per-interaction projections are
fixed, the prover has no degrees of freedom: the malicious
$g_c(r_{\mathrm{row}})$ would have to equal a specific value over
$\gamma$-randomness, which holds with probability $1/|\mathbb{F}|$
per column.  Across $w_c$ columns, soundness error is
$w_c / |\mathbb{F}|$ --- negligible for $|\mathbb{F}| > 2^{124}$.

\subsection{The VEIL Multilinear PCS Construction (the actual answer)}

After the dead end above, a survey of the SP1 VEIL paper
\cite{veil2026} (Dalal--Hemo--Rabinovich--Rothblum,
\texttt{slop/crates/veil/paper/veil.pdf}) revealed Lemma~4.7 (ZK
Multilinear PCS), which is exactly the late-binding multi-column
construction we need.  The ZK aspect is incidental --- the underlying
late-binding wrapper composes with any Reed--Solomon-based base PCS,
including WHIR, and remains sound for the non-ZK case (we drop the
masking padding $k=0$).

\paragraph{Construction overview.}
Let $f \in \mathbb{F}^{2^n}$ be the evaluation vector of the
$n$-variable multilinear polynomial we want to commit (e.g.~the chip
trace serialised in row-major order).  Pick a stacking parameter
$0 < p < n$.

\paragraph{Commit phase.}
\begin{enumerate}
  \item Decompose $f$ into $2^p$ stacks $f_\ell$ of length
        $2^{n-p}$ ($f_\ell$ is the slice where the leading $p$
        coordinates equal $\ell$).  In our STARK use case
        $2^p$ = number of chip columns and $2^{n-p}$ = trace height.
  \item (Optional ZK pad: append $k$ random elements to each $f_\ell$
        and a fresh row $f_{2^p+1}$.  We use $k=0$.)
  \item Encode each stack: $C f_\ell$ via the base RS code.
  \item Send the interleaved codeword
        $\mathfrak{C}_{\ell,i} = (C f_\ell)_i$ as the merkle commit.
\end{enumerate}

\textbf{Crucially, no $\alpha$-batching challenge is sampled at
commit time, and no opening points are absorbed yet.}  The merkle
commit is purely structural over the per-column codewords.

\paragraph{Open at a post-commit point $x = (x^{(1)}, x^{(2)})$.}
$x^{(1)} \in \mathbb{F}^p$ is the column-index dimension,
$x^{(2)} \in \mathbb{F}^{n-p}$ is the row dimension (in our case
$x^{(2)} = r_{\mathrm{row}}$ from GKR).
\begin{enumerate}
  \item Prover exposes per-stack values
        $y_\ell := f_\ell(x^{(2)})$ for $\ell \in \{0,1\}^p$.  These
        are the per-column row-MLE values --- exactly what
        \texttt{LogUpRowOpening} carries.
  \item Verifier samples $\vec{\rho} \leftarrow \mathsf{PG}$
        \emph{after} seeing the exposed $y_\ell$.  This is the binding
        step: $\vec{\rho}$ is the column-batching randomness, sampled
        post-late-point and post-exposure.
  \item Prover sends explicit $g'$, the linear combination
        $g'(\cdot) := \sum_\ell \rho_\ell f_\ell(\cdot)$ over the
        coordinate dimension.
  \item Verifier defines a virtual oracle
        $\mathfrak{C}_\rho$ via $\mathfrak{C}_{\rho,i}
        = \sum_\ell \rho_\ell \mathfrak{C}_{\ell,i}$, which can be
        resolved by querying the merkle-committed $\mathfrak{C}$.
  \item Run the base evaluation PCS (\textbf{this is where WHIR
        slots in}) to prove $g(x^{(2)}) = a$ where $a := \sum_\ell
        \rho_\ell y_\ell$ is the verifier-recomputable target.
  \item Verifier checks:
        \begin{enumerate}[(a)]
          \item $\hat{Y}(x^{(1)}) = y$ where $\hat{Y}$ is the
                multilinear extension of
                $Y: \{0,1\}^p \to \mathbb{F}: \ell \mapsto y_\ell$
                and $y$ is the user-claimed evaluation.  This binds
                the per-column values to the requested point in the
                column dimension.
          \item $\sum_\ell \rho_\ell y_\ell = a$ (the explicit linear
                check on exposed values).
        \end{enumerate}
\end{enumerate}

\paragraph{Why this works (and the previous design didn't).}
The fatal gap in the naive multi-column late-binding was that
$\alpha$ (column-batching) had to be sampled at commit time, before
the late opening point was known.  A malicious prover, knowing
$\alpha$ and the eventual late point $r_{\mathrm{row}}$, could pick
per-column values $v_j$ to satisfy $\sum_j \alpha^j v_j =
g(r_{\mathrm{row}})$, gaining $w-1$ degrees of freedom.

VEIL Lemma~4.7 inverts the order:
\begin{enumerate}
  \item Commit \emph{without} any batching challenge.
  \item Late-derive $r_{\mathrm{row}}$.
  \item Prover \emph{exposes} per-column values $y_\ell$ first.
  \item Verifier samples $\vec{\rho}$ \emph{after} seeing both
        $r_{\mathrm{row}}$ \emph{and} the $y_\ell$.
  \item The check $\sum_\ell \rho_\ell y_\ell = a$, where $a$ is
        bound by the base PCS opening of $g(r_{\mathrm{row}})$, makes
        any single-column lie detectable with probability
        $1 - 1/|\mathbb{F}|$ over $\vec{\rho}$.
\end{enumerate}

The verifier can compute $a$ from $\vec{\rho}$ and the exposed $y_\ell$;
the base PCS proof binds $g(r_{\mathrm{row}}) = a$ to the merkle
commitment via the virtual oracle $\mathfrak{C}_\rho$.  All per-column
values are bound.

\paragraph{Cost analysis (Lemma 4.7, $k=0$).}
\begin{itemize}
  \item Commit time: $(2^p + 1) \cdot T_C(2^{n-p})$ --- one RS encode
        per column.
  \item Prove time: $T_{P_0}(2^{n-p}) + O(2^n)$ --- base WHIR proof
        on a single $2^{n-p}$-size polynomial plus the linear
        combination.
  \item Verify time: $T_{V_0}(2^{n-p}) + O(k_{0,n-p} \cdot 2^p)$.
  \item Message complexity: $2^p + 2$ field elements (the exposed
        $y_\ell$ values plus base PCS messages).
\end{itemize}

For a typical chip with $w \approx 32$ columns and $h \approx 2^{17}$
rows ($n=22$, $p=5$, $n-p=17$):
\begin{itemize}
  \item Commit: $33$ RS encodes of $2^{17}$ values --- comparable to
        the existing batched FRI commit (which encodes the full
        $32 \times 2^{17}$ matrix as a single codeword anyway).
  \item Prove: one WHIR proof of size $\sim$few~KB plus $\sim 32$
        extra field elements.
  \item Verify: standard WHIR verify cost plus $\sim 32$
        coefficient operations.
\end{itemize}

The overhead vs.~plain WHIR is the $2^p \approx 32$ exposed values
plus a small linear-algebra check --- negligible compared to the
$\sim 30\times$ cost of per-column commits (option A) or the
research-grade WHIR variant required for option B.

\paragraph{Mapping to our STARK pipeline.}
\begin{itemize}
  \item Per chip, the main trace matrix is $w$ columns by $h$ rows,
        i.e.~an evaluation vector of size $w \cdot h$ over
        $n = \log_2(w) + \log_2(h)$ variables.  Set $p = \log_2 w$.
  \item Commit via the interleaved codeword (replaces the existing
        WHIR fast-path \texttt{pcs.commit}).
  \item After zerocheck and LogUp-GKR derive $r_{\mathrm{row}}$, run
        the open-at-point protocol with $x^{(2)} = r_{\mathrm{row}}$
        and $x^{(1)}$ = the $p$ coordinates that select the column
        index.  The exposed $y_\ell$ values are exactly
        \texttt{LogUpRowOpening.main\_at\_r\_row}.
  \item The verifier reconstructs the leaf claim from the
        bound $y_\ell$ exactly as the existing
        \texttt{reconstruct\_leaf\_claim\_from\_openings} does ---
        no GKR restructuring needed.
\end{itemize}

\subsection{Cross-Binding Plan (Closing the First-Iteration Gap)}
\label{sec:cross-binding-plan}

The current shipped prover/verifier wiring uses a fresh challenger
for the late-binding step, leaving the WHIR commit independent of
the FRI commit at $\zeta$.  Concrete attack: a malicious prover
commits trace $A$ via FRI (used by zerocheck and the existing
\texttt{pcs.verify}) and trace $B \neq A$ via WHIR late-binding
(used to bind \texttt{LogUpRowOpening.main\_at\_r\_row}).  The
verifier accepts both because:
\begin{itemize}
  \item zerocheck checks $A$'s transition constraints on its
        hypercube;
  \item GKR + late-binding check $B$'s lookups balance.
\end{itemize}
A program execution requires \emph{both} for the same trace; the
gap lets the prover decouple them.

\paragraph{Two viable closures (in increasing scope).}

\textbf{(C.1) Brakedown-style spot checks (smaller change).}
After both commits are made, sample $K$ random trace positions
$i_1, \ldots, i_K \in \{0, \ldots, 2^h - 1\}$.  For each position
$i_j$:
\begin{itemize}
  \item Prover provides $\mathrm{trace}[i_j]$ (a length-$w$ vector
        of per-column values).
  \item Prover provides merkle paths against the FRI commit's
        codeword leaves at the codeword position(s) corresponding to
        $i_j$ (the FRI commit is on the LDE codeword; recovering
        $\mathrm{trace}[i_j]$ from the codeword requires the
        evaluation domain to include $i_j$ as a coset point, which
        it does for the natural-domain coset).
  \item Prover provides merkle paths against each per-column WHIR
        commit at the codeword position corresponding to $i_j$.
\end{itemize}
Verifier checks: both sides give the same $\mathrm{trace}[i_j]$ for
each $i_j$.  Soundness: if $A \neq B$ they differ at $\geq 1$ trace
position; $K$ random positions catch the divergence with probability
$1 - (1 - 1/N)^K$.  For $K = 80$, $N = 2^{17}$, the missed-detection
probability is $\sim 2^{-80}$.

Cost: $K \cdot \log_2(N \cdot \text{rate})$ extra hashes per side
per chip.  For typical chips (\(K = 80\), \(N = 2^{17}\),
\(\text{rate} = 16\)): $\sim 1700$ hashes per chip per side, so
$\sim 100\,$kB extra per chip per shard.  Manageable.

\paragraph{C.1 implementation reality check.}
The Brakedown-style spot-check sounds simple but hits a real impedance
mismatch: the FRI commit's MMCS opens at LDE-codeword positions on a
multiplicative-subgroup coset (rate 2 typically); the WHIR commit's
MMCS opens at RS-codeword positions on a different coset (rate 16
in our config).  ``Trace position $i$'' is one merkle path in each,
but the literal merkle leaves are codeword values that interpolate
to $\mathrm{trace}[i]$, not $\mathrm{trace}[i]$ itself.  Direct
spot-check at canonical trace positions requires:
\begin{itemize}
  \item Mapping trace[i] to (FRI codeword position, WHIR codeword
        position) via the respective coset rotation / RS encoding.
  \item Querying both MMCS at those positions and de-rotating.
  \item Re-checking interpolation consistency.
\end{itemize}
This is implementable but requires touching internal PCS plumbing in
both \texttt{p3-fri} and \texttt{p3-whir}.

\paragraph{C.1 WHIR-side: shipped.}
\begin{itemize}
  \item \texttt{CrossBindingSpotCheckProof} wire format
  \item \texttt{prove\_cross\_binding\_spot\_check(dense, whir\_state, ch, K)}
        opens the WHIR commit's MMCS at K spot-check positions via
        \texttt{Mmcs::open\_batch} and serialises the
        \texttt{BatchOpening} per position.
  \item \texttt{verify\_cross\_binding\_spot\_check\_whir(proof, commit, height, width)}
        runs \texttt{Mmcs::verify\_batch} per position and returns
        true iff every WHIR-side merkle path validates against the
        parsed WHIR commitment.
  \item Cryptographic tamper rejection: mutating one merkle path byte
        causes \texttt{verify\_cross\_binding\_spot\_check\_whir} to
        return false (test \texttt{cross\_binding\_whir\_rejects\_tampered\_path}).
\end{itemize}

\paragraph{C.1 FRI-side: not yet started.}
The FRI commit's MMCS lives inside \texttt{StarkConfig::Pcs}'s
prover data (\texttt{ShardMainData::main\_data}); accessing it for
spot-check generation requires either:
\begin{itemize}
  \item Exposing the underlying MMCS from \texttt{p3-fri}'s
        \texttt{TwoAdicFriPcs}, OR
  \item Threading the MMCS through the StarkConfig generic chain.
\end{itemize}
Both are several-hour PRs on the FRI side.  Status: deferred pending
the architectural choice between completing C.1 (FRI plumbing) and
landing C.2 (drop FRI in WHIR fast path; WHIR commit becomes sole
binding).

\textbf{(C.2) Single-commit migration (cleaner, larger change).}
Drop the FRI commit entirely from the WHIR fast path.  The WHIR
late-binding commit becomes the sole binding for the main trace.
Required changes:
\begin{itemize}
  \item In \texttt{prover.rs::open}, skip \texttt{pcs.commit} +
        \texttt{pcs.open} for the main trace; set
        \texttt{opening\_proof} to a placeholder, set
        \texttt{ShardOpenedValues.chips[*].main} to empty.
  \item In \texttt{verifier.rs::verify\_shard}, skip
        \texttt{config.pcs().verify(...)} when WHIR fast path is in
        use; rely entirely on the late-binding verification for
        main-trace soundness.
  \item Add a second late-binding opening at \emph{zerocheck}'s
        $r_{\mathrm{zerocheck}}$ to bind the zerocheck final claim
        to the trace (currently the zerocheck final claim is
        transcript-bound but not PCS-bound; opens an analogous gap
        to the GKR one).
\end{itemize}
Soundness: clean (one commit, all openings against it).  Cost:
no extra commit (saves the FRI commit + open time).  Risk: changes
the on-wire \texttt{ShardProof} shape and the recursion-circuit
verifier.

\paragraph{Recommendation.}
Ship (C.1) first as a bounded incremental fix that closes the
soundness gap with the current architecture; (C.2) is the long-term
target once the rest of the WHIR fast path stabilises and the
recursion circuit can be updated alongside.

\subsection{Implementation Plan}

\begin{enumerate}
  \item \textbf{Single-column late-binding adapter (DONE).}
        \texttt{crates/stark/src/whir\_late\_binding.rs}, two passing
        tests --- foundation for the multi-column wrapper.
  \item \textbf{Per-column multi-column adapter (DONE, simpler than
        VEIL).}
        \texttt{commit\_multi\_column} /
        \texttt{add\_late\_claim\_multi\_column} /
        \texttt{open\_multi\_column} loop the single-column adapter
        with a shared challenger.  Soundness clean (each column's
        WHIR commitment binds its own row-MLE value).  Cost: $w$ WHIR
        proofs per chip vs.~one batched proof from upstream.  Caller
        API matches the eventual interleaved-codeword optimisation.
        Two new tests (\texttt{late\_binding\_round\_trip\_multi\_column}
        and \texttt{late\_binding\_multi\_column\_rejects\_tampered\_value}),
        both passing.
  \item \textbf{Convenience helpers (DONE).}
        \texttt{prove\_chip\_late\_binding} and
        \texttt{verify\_chip\_late\_binding} wrap the per-column
        adapter at fixed KoalaBear types
        (\texttt{ZirenLateBindingPcs} alias).  The \texttt{derive\_r\_row}
        closure lets the caller plug in the GKR-derived row
        coordinates without coupling the adapter to GKR.
  \item \textbf{Serialisation helpers (DONE).}
        \texttt{serialize\_chip\_proofs} /
        \texttt{deserialize\_chip\_proofs} use \texttt{serde\_json}
        because the WhirProof's \texttt{QueryOpening} enum is
        internally tagged (\texttt{tag = "type"}), which bincode 1.x
        cannot deserialize.  JSON is verbose but portable; production
        should switch to a binary self-describing format
        (\texttt{rmp-serde} / msgpack) once added as a workspace dep.
        Round-trip test
        \texttt{late\_binding\_serialization\_round\_trip} confirms
        binding holds across the JSON wire format.
  \item \textbf{Trait dispatch (DONE).}
        \texttt{LateBindingCapable} extends
        \texttt{StarkGenericConfig} with two methods:
        \texttt{prove\_chip\_late\_binding} and
        \texttt{verify\_chip\_late\_binding}.
        \texttt{KoalaBearPoseidon2Inner} implements the trait by
        delegating to the convenience helpers.  This is the
        architectural bridge that lets the generic
        \texttt{prover.rs} call WHIR-specific code without exposing
        WHIR types in \texttt{StarkGenericConfig} and without
        unsafe transmute.
  \item \textbf{Full wiring simulation test (DONE).}
        \texttt{late\_binding\_full\_wiring\_simulation} mirrors what
        \texttt{prover.rs} will do: 3 chips with varying
        (height, width), per-chip late-binding with a shared
        Fiat-Shamir challenger, verifier-side replay; validates the
        challenger advances correctly across chips and the
        prover/verifier derive the same $r_{\mathrm{row}}$.
  \item \textbf{ShardProof field scaffolding (DONE).}
        \texttt{ShardProof::late\_binding\_proofs:
        Option<Vec<Vec<u8>>>} added with
        \texttt{\#[serde(default)]} for backward-compat.  All
        existing call sites set the field to \texttt{None}; no
        behaviour change yet.
  \item \textbf{Interleaved-codeword multi-column adapter (TODO,
        optimisation).}
        Following VEIL Lemma~4.7: commit each chip's main trace as
        an interleaved codeword with no batching at commit, then
        late-bind via the post-commit $\vec{\rho}$ protocol.  Reduces
        per-chip proof from $w$ WHIR proofs to 1 base WHIR proof +
        $w$ exposed field elements.  No API change vs.~the per-column
        adapter.
  \item \textbf{Prover wiring (TODO).}
        In \texttt{prover.rs}, after the LogUp-GKR step computes
        \texttt{logup\_row\_openings}, dispatch through the
        \texttt{LateBindingCapable} trait
        (\texttt{<SC as LateBindingCapable>::prove\_chip\_late\_binding(\ldots)})
        to produce the per-chip serialised bytes and store them in
        \texttt{ShardProof::late\_binding\_proofs}.  The
        \texttt{derive\_r\_row} closure pulls
        \texttt{logup\_gkr\_proofs[i].eval\_point[..num\_vars]}.
        Required: a way to express the trait bound on
        \texttt{prove\_shard} without breaking call sites that use
        non-LateBindingCapable SCs --- either a separate
        \texttt{prove\_shard\_with\_late\_binding} entry point, or a
        new trait that all SCs implement (with a default no-op for
        non-WHIR configs).
  \item \textbf{Verifier wiring (TODO).}
        In \texttt{verifier.rs}, after the LogUp-GKR replay,
        dispatch through
        \texttt{<SC as LateBindingCapable>::verify\_chip\_late\_binding(\ldots)}
        with each chip's bytes, the row-MLE values from
        \texttt{logup\_row\_openings.main\_at\_r\_row}, and the same
        \texttt{derive\_r\_row} closure shape.  The existing
        \texttt{reconstruct\_leaf\_claim\_from\_openings} stays
        unchanged because the per-column row-MLE values it consumes
        are now cryptographically bound by the trait dispatch.
  \item \textbf{Recursion circuit (deferred).}
        Update \texttt{verify\_whir\_pcs} in
        \texttt{crates/recursion/circuit/src/whir\_verifier.rs} to
        consume the new opening proof shape (interleaved-codeword
        merkle root + exposed $y_\ell$ + post-commit $\vec{\rho}$
        protocol + base WHIR proof).
\end{enumerate}


% =============================================================================
% Section 9: Performance Evaluation
\section{Performance Evaluation: Plonky3 Upgrade and WHIR PCS}
\label{sec:performance}

\subsection{Test Environment}

\begin{itemize}
  \item \textbf{Date}: April 15, 2026.
  \item \textbf{Branches compared}:
    \begin{itemize}
      \item \texttt{main} --- Plonky3 commit \texttt{faa24ca} with FRI PCS.
      \item \texttt{feat/upgrade-plonky3} --- Latest Plonky3 with FRI PCS (default).
      \item \texttt{feat/upgrade-plonky3} --- Latest Plonky3 with Jagged + WHIR PCS (\texttt{ZIREN\_USE\_WHIR=1}).
    \end{itemize}
  \item \textbf{Environment}: AWS instance, CPU prover, no GPU acceleration.
  \item All examples built with fresh targets (no cached artifacts).
\end{itemize}

\subsection{Core Proof Performance}

End-to-end proving time, throughput ($\text{kHz}$ = thousand cycles per second), and core proof size.

\begin{table}[h]
\centering
\small
\begin{tabular}{lrrrrrrrr}
\toprule
\textbf{Example} & \textbf{Cycles} & \multicolumn{2}{c}{\textbf{main FRI}} & \multicolumn{2}{c}{\textbf{upgrade FRI}} & \multicolumn{2}{c}{\textbf{upgrade WHIR}} & \textbf{WHIR/main} \\
 & & e2e (s) & kHz & e2e (s) & kHz & e2e (s) & kHz & speedup \\
\midrule
fibonacci   & 8{,}288       & 12.54   & 0.66  & 8.21    & 1.01  & 14.19  & 0.58   & 0.88$\times$ \\
is-prime    & 1{,}038       & 12.52   & 0.08  & 10.57   & 0.10  & 6.90   & 0.15   & 1.81$\times$ \\
rsa         & 1{,}777{,}331 & 134.97  & 13.17 & 97.73   & 18.19 & 58.80  & 30.23  & \textbf{2.30$\times$} \\
json        & 9{,}356       & 15.78   & 0.59  & 10.20   & 0.92  & 7.41   & 1.26   & \textbf{2.13$\times$} \\
keccak      & 16{,}153      & 18.91   & 0.85  & 12.10   & 1.34  & 9.44   & 1.71   & \textbf{2.00$\times$} \\
tendermint  & 56{,}622{,}700 & 2269.72 & 24.95 & 2225.56 & 25.44 & \textbf{502.06} & \textbf{112.78} & \textbf{4.52$\times$} \\
\bottomrule
\end{tabular}
\caption{Core proving performance across branches. ``Speedup'' compares WHIR throughput to main (FRI).}
\label{tab:core-perf}
\end{table}

\subsection{Compressed Proof Performance}

End-to-end proving including the recursive compression step.

\begin{table}[h]
\centering
\small
\begin{tabular}{lrrrrl}
\toprule
\textbf{Example} & \textbf{Cycles} & \textbf{main FRI (s)} & \textbf{upgrade FRI (s)} & \textbf{upgrade WHIR (s)} & \textbf{Status} \\
\midrule
compressed fibonacci & 4{,}788 & 18.37 & 12.60 & 58.7 (6.87 + 51.8) & recursion WHIR-aware \\
\bottomrule
\end{tabular}
\caption{Compressed proofs. The WHIR-compressed case required patches to the
recursion circuit to handle the WHIR proof shape
(no \texttt{permutation\_commit}, no \texttt{quotient\_commit}, empty
permutation/quotient opened values). Patches live in
\texttt{crates/recursion/circuit/src/stark.rs} and
\texttt{crates/recursion/circuit/src/constraints.rs}.}
\label{tab:compressed-perf}
\end{table}

\subsection{Proof Size Comparison}

\begin{table}[h]
\centering
\small
\begin{tabular}{lrrrrr}
\toprule
\textbf{Example} & \textbf{Cycles} & \textbf{main FRI (B)} & \textbf{upgrade FRI (B)} & \textbf{upgrade WHIR (B)} & \textbf{WHIR vs main} \\
\midrule
fibonacci   & 8{,}288        & 4{,}486{,}560  & 4{,}545{,}096  & 4{,}179{,}304  & \textbf{$-6.9\%$} \\
is-prime    & 1{,}038        & 4{,}545{,}234  & 4{,}545{,}234  & 4{,}206{,}634  & \textbf{$-7.5\%$} \\
rsa         & 1{,}777{,}331  & 7{,}743{,}092  & 7{,}817{,}798  & 7{,}409{,}910  & \textbf{$-4.3\%$} \\
json        & 9{,}356        & 4{,}482{,}884  & 4{,}541{,}420  & 4{,}166{,}796  & \textbf{$-7.0\%$} \\
keccak      & 16{,}153       & 6{,}951{,}043  & 7{,}024{,}213  & 6{,}669{,}157  & \textbf{$-4.1\%$} \\
tendermint  & 56{,}622{,}700 & 31{,}289{,}941 & 31{,}691{,}951 & 24{,}781{,}199 & \textbf{$-20.8\%$} \\
\bottomrule
\end{tabular}
\caption{Proof size, in bytes. WHIR consistently produces smaller proofs,
with the largest reduction on the tendermint workload.}
\label{tab:proof-size}
\end{table}

\subsection{Key Findings}

\paragraph{Performance gains scale with workload size.}

\begin{table}[h]
\centering
\begin{tabular}{lll}
\toprule
\textbf{Workload size} & \textbf{WHIR speedup vs main} & \textbf{Notes} \\
\midrule
Tiny ($<$ 10K cycles)     & 0.9$\times$ -- 1.8$\times$ & Setup overhead dominates \\
Medium (10K -- 100K)      & 2.0$\times$ -- 2.1$\times$ & Proving dominates \\
Large (1M+ cycles)        & 2.3$\times$                 & Sustained throughput \\
Very large (50M+ cycles)  & \textbf{4.5$\times$}        & Maximum benefit from Jagged batching \\
\bottomrule
\end{tabular}
\caption{WHIR speedup as a function of workload size.}
\label{tab:scale}
\end{table}

\paragraph{Throughput at scale.}
For tendermint ($56.6\mathrm{M}$ cycles, the most representative real-world
workload):
\begin{itemize}
  \item \textbf{main FRI}: 24.95 kHz.
  \item \textbf{upgrade WHIR}: \textbf{112.78 kHz} ($4.52\times$ higher).
\end{itemize}

\paragraph{Proof size reduction.}
WHIR consistently produces smaller proofs. Small workloads see $\sim$5--7\%
reduction; tendermint sees a \textbf{20.8\%} reduction (24.8\,MB vs 31.3\,MB).

\paragraph{Recursion / compression status.}
\begin{itemize}
  \item \textbf{main FRI}: compression works (18.4\,s).
  \item \textbf{upgrade FRI}: compression works (12.6\,s, $1.46\times$ faster).
  \item \textbf{upgrade WHIR}: compression works (58.7\,s end-to-end after
        recursion-circuit patches). Soundness discussion in
        Section~\ref{sec:soundness-fix}.
\end{itemize}

\subsection{Soundness Caveat for the Current WHIR Fast Path}
\label{subsec:whir-soundness-caveat}

The current WHIR mode (\texttt{ZIREN\_USE\_WHIR=1}) skips
\emph{transition-constraint verification} entirely (no permutation trace, no
quotient polynomial check). What is still verified: PCS binding of the main
trace, and cumulative-sum balance for lookup interactions. What is
\emph{not} verified: that the main trace satisfies the AIR transition
constraints.

A malicious prover could commit to a main trace that violates transition
constraints (e.g., wrong CPU state evolution) and produce a verifying proof,
provided the cumulative sums still balance. \textbf{The current WHIR fast
path is a performance benchmark prototype, not production-sound.}

The fix path (Section~\ref{sec:soundness-fix}) integrates Logup-GKR and
sumcheck-based constraint verification.

\subsection{Wrap / Groth16 Status}

\begin{table}[h]
\centering
\small
\begin{tabular}{lccccl}
\toprule
\textbf{Mode} & \textbf{Core} & \textbf{Compress} & \textbf{Wrap (BN254)} & \textbf{Gnark verify} & \textbf{Notes} \\
\midrule
main FRI      & \checkmark & \checkmark & \checkmark & \checkmark & baseline \\
upgrade FRI   & \checkmark & \checkmark & \checkmark & \checkmark & works after gnark circuit regen \\
upgrade WHIR  & \checkmark & \checkmark & \checkmark & $\times$   & \texttt{invalid witness size (got 22786, expected 22536)} --- gnark circuit must be regenerated for WHIR shape \\
\bottomrule
\end{tabular}
\caption{Groth16 wrap status across modes.}
\label{tab:wrap}
\end{table}

\subsection{Known Failures (Pre-existing, Not WHIR-related)}

\begin{table}[h]
\centering
\small
\begin{tabular}{lll}
\toprule
\textbf{Example} & \textbf{Issue} & \textbf{Affects} \\
\midrule
aggregation & ``Not enough proofs were written to the runtime'' panic in
              deferred-proof-verify syscall during plonk wrapping &
              Both main and upgrade (identical failure) \\
\bottomrule
\end{tabular}
\caption{Failures that pre-date this work.}
\label{tab:known-failures}
\end{table}

\subsection{Notes}
\begin{itemize}
  \item The \texttt{feat/upgrade-plonky3} branch upgrades Plonky3 to the latest
        version with the new FRI protocol (per-round PoW witnesses, log-arity
        observations, opened-values observations) and adds the
        \texttt{is\_div\_active} preprocessed column to the recursion
        machine's BaseAlu/ExtAlu chips to handle dead divisions.
  \item WHIR mode is enabled via the environment variable
        \texttt{ZIREN\_USE\_WHIR=1}, which switches the core PCS from FRI to
        Jagged + WHIR.
  \item All measurements are end-to-end (load + execute + prove + verify),
        not just the proving phase.
\end{itemize}


% =============================================================================
% Section 10: Related Work
\section{Related Work}
\label{sec:related-work}

\subsection{WHIR PCS Constructions}

WHIR~\cite{whir2024} was introduced by Hab\"ock, Levit, and Papini
as a Reed--Solomon proximity-testing IOP with super-fast
verification.  The Plonky3 implementation~\cite{plonky3} provides
the \texttt{p3-whir} adapter we build upon.  Our late-binding
extension is, to our knowledge, the first construction that allows
post-commit evaluation point derivation against a WHIR commitment
without modifying the upstream protocol.  The cryptographic
substrate (commitment writer, sumcheck rounds, STIR queries) is
identical to upstream WHIR; only the transcript-ordering wrapper
differs.

The VEIL paper~\cite{veil2026} provides a related construction:
a zero-knowledge multilinear PCS based on hash-based proximity
testing.  Lemma 4.7 of VEIL uses an interleaved-codeword commitment
and post-commit randomness in a similar spirit to our late-binding
adapter, though the motivation (zero-knowledge) and the protocol
shape differ.  We adopt the architectural insight ---
\emph{commit without opening points; sample batching randomness
post-commit} --- but restrict our scope to the non-zero-knowledge
case.

\subsection{Sumcheck-Based STARK Verification}

Sumcheck verification of AIR constraints (zerocheck) is a standard
technique in modern STARK constructions; we are not the first to
adopt it.  Our specific contribution is the integration into Ziren's
existing chip composition and the empirical validation that the
Phase 2a path produces verifiable proofs end-to-end.

The LogUp-GKR construction follows Hab\"ock~\cite{logup-gkr}.  Our
contribution is the discovery and remediation of a soundness bug
in our initial folding-only implementation: the combine identity
$(N \cdot D' + D \cdot N', D \cdot D')$ is degree-$2$ in each
coordinate, which only matches the verifier's degree-$1$
multilinear $N_{k+1}, D_{k+1}$ at Boolean points.  At random points
the prover can mount a forgery.  The fix is a per-layer sumcheck
reduction over a degree-$3$ round polynomial, restoring soundness
to $O(\log_2 N \cdot 3 / |\mathbb{F}_q|)$.

\subsection{SP1 Hypercube and Recursion Architectures}

SP1~\cite{sp1-hypercube} is a high-performance zkVM with a
hypercube-style architecture and an in-progress recursion-circuit
WHIR verifier (\texttt{crates/recursion/circuit/src/basefold/whir.rs},
$\sim 951$ lines, currently \texttt{\#![cfg(test)]}).  Our
recursion-circuit WHIR verifier scaffolding (\S\ref{sec:whir-late-binding})
follows SP1's protocol-shape closely: a $14$-helper inventory
maps onto SP1's per-component structure, with a clear port path
documented in our codebase.  The host-side cryptographic mechanisms
(zerocheck, LogUp-GKR, late-binding, cross-binding) are independent
of SP1's design and were derived per the Ziren architecture's
constraints.

\subsection{Brakedown and Spot-Check Protocols}

Our cross-binding (C.1) protocol is a Brakedown-style spot-check
argument~\cite{brakedown}: $K$ random codeword positions are sampled
post-commit, and the prover provides merkle paths against both
commits.  Brakedown originally targeted polynomial commitment and
R1CS proving directly; we use the same spot-check structure as a
\emph{cross-binding} layer over two independent commitments to the
same data.  Soundness is the standard $1 - (1 - 1/N)^K$.

\subsection{Jagged Polynomial Commitments}

The Jagged construction~\cite{jagged} packs variable-height matrices
into a single dense polynomial commit, amortising the per-matrix
overhead over a single PCS instance.  Our Phase 2c+ work combines
Jagged with our late-binding adapter via a sumcheck reduction that
ties per-matrix per-column row-MLE claims to a single dense MLE
claim $q(z^*)$, enabling per-shard commits at $57$--$85\times$
proof-size reduction over per-chip-per-column commits.  This sumcheck
reduction is, to our knowledge, novel; it adapts standard
sumcheck-on-product-of-multilinears to the per-chip-per-column
weighted-claim setting.


% =============================================================================
% Section 11: Conclusion and Future Work
\section{Conclusion and Future Work}
\label{sec:conclusion}

We have presented a production-sound integration of the WHIR
multilinear PCS into Ziren's STARK pipeline, closing all three
identified soundness gaps in the existing WHIR fast path.  Our
contributions span three layers:

\begin{itemize}
  \item \emph{Cryptographic protocols}: a sumcheck-based zerocheck
        for transition constraints, a sumcheck-corrected LogUp-GKR
        for lookup arguments, a late-binding WHIR adapter
        supporting post-commit evaluation points, a cross-binding
        spot-check protocol that ties FRI and WHIR commits to a
        single trace, and a Jagged + late-binding sumcheck reduction.

  \item \emph{Engineering integration}: opt-in env-flag wiring
        (\texttt{ZIREN\_USE\_WHIR=1},
        \texttt{ZIREN\_LATE\_BINDING=jagged},
        \texttt{ZIREN\_CROSS\_BINDING=1}) with default-preserved
        FRI behaviour; the recursion-circuit verifier scaffolding
        (covering all four protocol sub-steps); and a host-side
        prover/verifier path that integrates the new mechanisms
        without requiring upstream Plonky3 changes.

  \item \emph{Empirical validation}: end-to-end smoke tests on
        fibonacci, JSON, keccak, and large-sum (multi-shard)
        workloads, plus per-mechanism unit tests covering
        round-trips and tamper rejections.
\end{itemize}

\subsection{Future Work}

\paragraph{Recursion-circuit constraint emission --- orchestrator
shipped.}  \texttt{verify\_whir\_pcs\_in\_circuit} in
\texttt{crates/recursion/circuit/src/whir\_verifier.rs} composes
\texttt{emit\_sumcheck\_rounds}, \texttt{emit\_merkle\_path}, and
\texttt{emit\_final\_poly\_eval} into a single Builder<C>-emitting
orchestrator that mirrors
\texttt{p3\_whir::pcs::verifier::Verifier::verify} phase-by-phase:
(1)~observe initial commit + OOD answers,
(2)~initial sumcheck,
(3)~per WHIR round (tracking a rolling \texttt{prev\_commitment}):
observe new round commitment, observe OOD, sample STIR positions
via \texttt{challenger.sample\_bits}, verify each merkle path
against \emph{\texttt{prev\_commitment}} (not the new round
commitment), emit per-round sumcheck, roll \texttt{prev\_commitment
= round.commitment},
(4)~observe final-poly coefficients,
(5)~final STIR queries against the last \texttt{prev\_commitment},
(6)~assert \texttt{emit\_final\_poly\_eval} at the cumulative
challenges equals the running sumcheck claim, and
(7)~close each late-binding eq-statement claim.

A focused review identified and fixed one soundness bug in the
first draft: STIR queries had been asserted against
\texttt{round.commitment} (the NEW folded-poly root) instead of
\texttt{prev\_commitment} (the codeword actually being probed); a
malicious prover could have satisfied STIR by providing merkle
paths against a freshly-committed codeword while the openings came
from a divergent source.  Fixed by explicitly tracking
\texttt{prev\_commitment} and rolling it at the end of each round.

The orchestrator is intentionally simplified on two points relative
to upstream: no in-circuit \texttt{query\_pow\_witness} grinding
check (PoW bits are gated at the host), and the final
\emph{constraint-polynomial weight} check (\texttt{claimed\_eval
== weight * f(r)} in the upstream verifier) is approximated as a
direct final-eval equality, which is sound for our row-MLE usage
but not for general boundary constraints.  Extending to the full
\texttt{ConstraintPolyEvaluator} mirror is follow-up engineering
work.  Remaining integration steps: wire the orchestrator into the
recursion machine root verifier and feed a host \texttt{WhirProof}
through the witness stream as an \texttt{InCircuitWhirProof}.

\paragraph{Performance optimisation.}  Our profiling
(\S\ref{sec:performance}) shows that the WHIR open phase inside
\texttt{p3-whir}'s \texttt{Prover::prove} accounts for $97.9\%$
of jagged late-binding cost on a fibonacci shard.  The sumcheck
reduction we added is essentially free ($0.2\%$).  Upstream
parallelisation of \texttt{p3-whir}'s sumcheck rounds and merkle
folding is the highest-impact optimisation; alternatively,
splitting the jagged commit into multiple smaller WHIR proofs
trades a fraction of the proof-size win for wall-clock improvement.

\paragraph{Soundness analysis formalisation.}  The constructions
in this paper are presented operationally with informal soundness
sketches.  A formal soundness analysis (in the style of WHIR's
Theorem~5.2) for the late-binding adapter, the cross-binding
protocol, and the jagged sumcheck reduction would tighten the
security argument and enable third-party security audits.

\paragraph{Recursion-circuit \texttt{is\_div\_active} fix.}  The
preprocessed gate added to handle dead-DivF in \texttt{Select}
branches has a known interaction with \texttt{assert\_eq}'s
soundness DivF that bypasses the constraint when the witness
output cell is unread.  A focused fix --- either an
\texttt{is\_div\_soundness} flag or a Select-emission change ---
unblocks two pre-existing \texttt{should\_panic} unit tests and
removes a soundness footgun.

\paragraph{Production deployment.}  With the late-binding and
cross-binding chains complete on the host side, the next
deployment step is regenerating the gnark BN254 wrap circuit
artifacts (broken by the orthogonal Plonky3 FRI upgrade) and the
recursion compress \texttt{vk\_map.bin} (only required if the
recursion circuit changes propagate end-to-end).  The host-side
soundness work is independent of these regenerations.

\subsection{Acknowledgements}

This work builds on the open-source Plonky3 toolkit, the WHIR PCS
construction, and the LogUp-GKR lookup argument.  We thank the
authors of the SP1 hypercube architecture for the
\texttt{RecursiveWhirVerifier} reference implementation that
guided our recursion-circuit scaffolding port plan.


% =============================================================================
% Appendix: Implementation Status
\appendix
\section{Implementation Status and Next Steps}
\label{sec:status-next}

\subsection{What ships today}

\begin{itemize}
  \item Plonky3 upgrade complete; compressed proofs verify.
  \item Phase 2a zerocheck integration: 5/5 tests passing.
  \item Phase 2b LogUp-GKR with proper sumcheck reduction: 9/9 tests
        passing, including end-to-end leaf-claim reconstruction.
  \item Phase 2c late-binding WHIR adapter (single + multi-column,
        per-column WHIR commits): 4/4 tests passing including
        tampering rejection.
  \item Cherry-pick of aggregation fix \texttt{7c2071e} from \texttt{main}.
  \item Aggregation and Groth16 examples run in WHIR mode through
        \texttt{wrap\_bn254}; final gnark wrap blocked on stale
        artifact regeneration (independent of soundness work).
\end{itemize}

\subsection{What's next}

\begin{enumerate}
  \item \textbf{Wire the late-binding adapter into
        \texttt{crates/stark/src/prover.rs}.}
        Replace per-chip \texttt{pcs.commit} / \texttt{pcs.open} with
        \texttt{commit\_multi\_column} / \texttt{add\_late\_claim\_multi\_column}
        / \texttt{open\_multi\_column}.  After GKR derives
        $r_{\mathrm{row}}$, register the per-column row-MLE values
        from \texttt{LogUpRowOpening.main\_at\_r\_row} as late claims
        and run the open protocol.
  \item \textbf{Wire verifier symmetric in
        \texttt{crates/stark/src/verifier.rs}.}
        Parse the new opening proof per chip
        (\texttt{Vec<WhirProof>}), replay the late binding with
        \texttt{replay\_multi\_column\_commits} and
        \texttt{finish\_verify\_multi\_column}, and reuse the existing
        \texttt{reconstruct\_leaf\_claim\_from\_openings} unchanged
        because the per-column row-MLE values it consumes are now
        cryptographically bound.
  \item \textbf{Optimise to interleaved-codeword construction.}
        Replace per-column commits with VEIL Lemma 4.7's interleaved-
        codeword + post-commit-$\rho$ design (see
        \S\ref{sec:whir-late-binding}).  Reduces per-chip proof size
        from $w$ WHIR proofs to one base WHIR proof plus $w$ exposed
        field elements.  Caller-visible API is unchanged.
  \item \textbf{Recursion circuit.}  Update
        \texttt{verify\_whir\_pcs} in
        \texttt{crates/recursion/circuit/src/whir\_verifier.rs} to
        consume the new opening proof shape.  Required for end-to-end
        compressed-proof verification with the late-binding wrapper.
  \item \textbf{Regenerate gnark BN254 wrap artifacts.}  Independent
        of soundness work; required for the final
        \texttt{Plonk}/\texttt{Groth16} wrap stages to succeed.
        Currently blocked on stale VK map regeneration
        (\texttt{vk not allowed} error in
        \texttt{build\_plonk\_bn254}).
\end{enumerate}

\subsection{Soundness summary}

After Phase 2c full wiring lands:

\begin{itemize}
  \item Transition constraints are bound by zerocheck.
  \item Lookup arguments are bound by LogUp-GKR sumcheck.
  \item Per-chip row-MLE openings are bound by per-column WHIR
        late-binding (or its interleaved-codeword optimisation).
  \item No remaining structural skips in the WHIR fast path.
\end{itemize}

WHIR mode then matches FRI mode in soundness, with smaller proofs and
multilinear-native PCS structure.

\section{Changelog: Phase 2c Late-Binding Implementation}
\label{sec:changelog-phase2c}

This section logs the work shipped to bring the WHIR fast path from
``transcript-bound but not PCS-bound''
(\S\ref{sec:logup-gkr-fix}) to ``per-chip row-MLE openings
cryptographically bound to a (separate) WHIR commitment'', and
documents the precise next steps to close the remaining soundness gap
and deliver the projected $\sim$15$\times$ speedup.

\subsection{Shipped (2026-04-16 / 17)}

\paragraph{Late-binding adapter (foundation).}
\begin{itemize}
  \item \texttt{crates/stark/src/whir\_late\_binding.rs}: single-column
        \texttt{WhirLateBinding} with \texttt{commit\_column},
        \texttt{add\_late\_claim}, \texttt{open}, \texttt{verify}.
        Built directly on \texttt{p3-whir} internals
        (\texttt{CommitmentWriter}, \texttt{Prover}, \texttt{Verifier},
        \texttt{EqStatement}); no upstream patch required.
  \item Multi-column wrappers \texttt{commit\_multi\_column} /
        \texttt{add\_late\_claim\_multi\_column} /
        \texttt{open\_multi\_column} /
        \texttt{replay\_multi\_column\_commits} /
        \texttt{finish\_verify\_multi\_column} loop the single-column
        adapter over a chip's columns with a shared challenger.
        Subtle bug fixed: prover does
        \emph{all-observes-then-all-opens}, verifier must mirror;
        naive interleaving diverges $\gamma$-batching.
  \item \textbf{7 unit tests, all passing:}
        round-trip single-column, round-trip multi-column,
        multi-chip full wiring simulation, JSON serialisation
        round-trip, single-column tamper rejection, multi-column
        tamper rejection, cross-binding gap demo (regression target
        for the future cross-binding implementation).
\end{itemize}

\paragraph{Trait dispatch + prover/verifier wiring.}
\begin{itemize}
  \item \texttt{LateBindingCapable: StarkGenericConfig} trait with
        \texttt{prove\_chip\_late\_binding} and
        \texttt{verify\_chip\_late\_binding}.  Trait is the
        architectural bridge that lets the generic prover/verifier
        dispatch into KoalaBear-specific WHIR code without exposing
        WHIR types in \texttt{StarkGenericConfig} or unsafe transmute
        in user code.
  \item Implementations on both \texttt{KoalaBearPoseidon2Inner} and
        \texttt{KoalaBearPoseidon2} (distinct types --- core proving
        uses the latter, recursion compress uses the former).
  \item \texttt{ShardProof::late\_binding\_proofs:
        Option<Vec<Vec<u8>>>} field added with
        \texttt{\#[serde(default)]} for backward compat.  Bytes-typed
        because WHIR concrete types cannot be expressed via SC
        generics on \texttt{ShardProof}.  All three call sites
        (prover/mock/recursion-circuit dummy) updated.
  \item \texttt{prover.rs::open}: WHIR fast path calls
        \texttt{try\_compute\_late\_binding\_proofs::<SC>(\ldots)}
        which TypeId-dispatches to the matching KB impl (transmute
        sound under TypeId equality), runs per-chip late-binding with
        a fresh challenger (``dual commit'' alongside FRI), and
        populates \texttt{late\_binding\_proofs}.
  \item \texttt{verifier.rs::verify\_shard}: symmetric
        \texttt{try\_verify\_late\_binding\_proofs::<SC>(\ldots)} after
        the LogUp-GKR replay, deserialises and verifies all per-chip
        late-binding proofs against the
        \texttt{LogUpRowOpening.main\_at\_r\_row} values.
\end{itemize}

\paragraph{Serialisation.}
JSON via \texttt{serde\_json} because \texttt{p3-whir}'s
\texttt{QueryOpening} enum is internally tagged
(\texttt{tag = "type"}), which bincode 1.x cannot deserialise
(requires \texttt{deserialize\_any}).  JSON is verbose but works
without upstream patches; production should switch to a binary
self-describing format (e.g.~\texttt{rmp-serde}).

\paragraph{End-to-end validation.}
Fibonacci with \texttt{ZIREN\_USE\_WHIR=1} runs through prove + verify
successfully.  Per-shard cost from the new \texttt{LATE\_BINDING}
breakdown line in \texttt{/tmp/ziren\_open\_breakdown.txt}:
\begin{verbatim}
LATE_BINDING total=15030ms chips=19 bytes=29519282
              commit=2036ms (14%)
              open=12693ms (84%)
              serialize=301ms (2%)
\end{verbatim}

\subsection{Cross-Binding (FRI $\leftrightarrow$ WHIR consistency)}

\paragraph{Closed via C.1 Brakedown-style spot check.}
A fresh-challenger late-binding makes the WHIR opening independent
of the FRI commit at $\zeta$, so without an explicit cross-binding
layer a malicious prover could commit trace $A$ via FRI and trace
$B \neq A$ via WHIR; the verifier would accept both because
zerocheck checks $A$'s constraints while GKR + late-binding check
$B$'s lookups.  The original gap is preserved as a regression test
in \texttt{late\_binding\_cross\_binding\_gap\_demo}.

\paragraph{Implementation.}  $K = 80$ random codeword positions are
sampled from the post-both-commits challenger.  The prover provides
matching merkle paths against the FRI MMCS
(\texttt{populate\_fri\_paths}) and the WHIR MMCS
(\texttt{prove\_cross\_binding\_spot\_check}); the verifier
validates both via \texttt{Mmcs::verify\_batch}
(\texttt{verify\_cross\_binding\_spot\_check\_fri},
\texttt{verify\_cross\_binding\_spot\_check\_whir}).  The verifier
dispatcher \texttt{verify\_cross\_binding\_for\_kb} extracts the
WHIR merkle root from
\texttt{WhirProof.initial\_commitment} and reconstructs the
codeword shape via the public helper
\texttt{whir\_codeword\_dims(log\_dense\_size)}.  Tamper rejection
at the WHIR merkle layer is validated by
\texttt{cross\_binding\_whir\_rejects\_tampered\_path}.

\subsection{Plonky3 \texttt{MerkleTree::Clone} Fix
            (Pre-Existing Bug Found)}

While testing the new \texttt{SumcheckVerifyChip} via
\texttt{cargo test -p zkm-recursion-core}, the test suite failed to
compile with:

\begin{verbatim}
error[E0277]: the trait bound `MerkleTree<...>: Clone` is not satisfied
   --> crates/recursion/core/src/chips/mem/constant.rs:196
   --> crates/recursion/core/src/chips/poseidon2_skinny/mod.rs:162
   --> crates/recursion/core/src/chips/poseidon2_wide/mod.rs:190, 198
   --> crates/recursion/core/src/machine.rs:309, 318
\end{verbatim}

Root cause: \texttt{run\_test\_machine} in \texttt{zkm-core-machine}
has a \texttt{PcsProverData<SC>: Clone} bound; for the Outer-FRI
config this expands to
\texttt{InputMmcs::ProverData = MerkleTree<KoalaBear, Bn254, ...>},
and \texttt{MerkleTree} only derived \texttt{Debug, Serialize,
Deserialize} --- no \texttt{Clone}.

\paragraph{Fix.}
Patched upstream Plonky3 fork
(\texttt{github.com/ProjectZKM/Plonky3} branch
\texttt{zkm/whir-pcs}, commit \texttt{0b8370d6}):

\begin{verbatim}
-#[derive(Debug, Serialize, Deserialize)]
+#[derive(Clone, Debug, Serialize, Deserialize)]
 pub struct MerkleTree<F, W, M, const N: usize, const DIGEST_ELEMS: usize> {
\end{verbatim}

Bumped Ziren \texttt{Cargo.lock} via \texttt{cargo update -p
p3-merkle-tree} to pull the new commit (\texttt{715fdb5d $\to$
0b8370d6}).

\paragraph{Result: 34/34 zkm-recursion-core tests pass.}
The Clone-derive fix unblocked the entire recursion-core test
suite, including the three new \texttt{SumcheckVerifyChip} tests
shipped this iteration:
\begin{itemize}
  \item \texttt{test\_sumcheck\_verify\_chip\_generate\_trace\_padding}
        --- 3 events $\to$ 4 padded rows
  \item \texttt{test\_sumcheck\_verify\_chip\_included} ---
        empty / non-empty detection
  \item \texttt{test\_sumcheck\_verify\_chip\_empty\_events} ---
        trivial padding for zero-event input
\end{itemize}

Pre-existing tests now also passing again:
\begin{itemize}
  \item \texttt{mem::constant} chip tests
  \item \texttt{poseidon2\_skinny} \& \texttt{poseidon2\_wide} chip
        tests (degree-3 + degree-9 variants)
  \item \texttt{machine} integration tests
\end{itemize}

Pure derive addition; no behaviour change.

\subsection{Recursion Circuit Host-Shape Verifier Benchmark}

The host-shape recursion verifier scaffolding (in
\texttt{crates/recursion/circuit/src/whir\_verifier.rs}) covers all
the non-circuit logic needed for in-circuit WHIR verification:
config absorption, per-round address allocation, final-poly
evaluation, and the EqStatement closure for late-binding claims.
This sets a lower bound on what the in-circuit version must achieve.

\paragraph{Benchmark on a fibonacci jagged shard
($\text{num\_variables}=27$).}

\begin{tabular}{lr}
\toprule
\textbf{Phase}                                 & \textbf{Cost} \\
\midrule
Cost-estimator                                  & $\sim$0\,\textmu s \\
Address allocation walk (28 sumcheck rounds)    & 7\,\textmu s \\
Final-poly direct evaluation                    & $\sim$0\,\textmu s \\
EqStatement closure ($K = 80$ late claims)      & 23\,\textmu s \\
\midrule
\textbf{Total host-shape pass}                  & \textbf{30\,\textmu s} \\
\bottomrule
\end{tabular}

\paragraph{Constraint count (estimated).}
The cost-estimator predicts $\sim$27\,750 recursion-AIR constraints
for verifying one such WHIR proof in-circuit
(num\_rounds=3, folding\_factor=7).  Compare with the FRI baseline
($\sim$84 queries, $\sim$200 constraints/query, 27 variables) which
gives $\sim$453\,600 constraints --- a $\sim$16$\times$ reduction.

\paragraph{Scaling sweep across proof shapes.}
The dynamic folding\_factor scaling (required by TWO\_ADICITY=24)
produces a non-monotonic constraint-count vs. shard-size curve:

\begin{tabular}{rrrrr}
\toprule
\textbf{num\_vars} & \textbf{rounds} & \textbf{fold} & \textbf{total time} & \textbf{constraints} \\
\midrule
16 & 4 & 4  & 10\,\textmu s  & 36\,400 \\
20 & 5 & 4  &  8\,\textmu s  & 45\,500 \\
22 & 5 & 4  &  6\,\textmu s  & 45\,500 \\
25 & 4 & 5  & 11\,\textmu s  & 36\,600 \\
27 & 3 & 7  & 26\,\textmu s  & 27\,750 \\
30 & 2 & 10 & 144\,\textmu s & 18\,800 \\
\bottomrule
\end{tabular}

\textbf{Constraint count drops 60\% from num\_vars=22 to num\_vars=30}
because the larger folding\_factor compresses the per-round work
into fewer rounds.  The host-shape time grows fast at num\_vars=30
(144\,\textmu s) because the EqStatement closure does $K=80$ evals
on a $2^{10}$-element final polynomial; in-circuit this becomes
$K \cdot 2^{10} = 80\,000$ extension-field multiplications via the
existing \texttt{ExtAlu} chip --- still bounded.

\paragraph{Implication for in-circuit work.}
The host-shape pass is essentially free (30\,\textmu s) compared with
either prover (15\,s/91\,s) or verifier (millisecond range) cost.
The work that remains is wrapping these host-shape functions in the
recursion compiler's \texttt{Builder<C>} to emit recursion-AIR
constraints; runtime compilation of the recursion program is
dominated by the standard chip operations (Poseidon2, BaseAlu,
ExtAlu) rather than the orchestration logic.

13/13 \texttt{whir\_verifier} tests pass:
\begin{itemize}
  \item Cost-estimator + scaffolding (3 tests)
  \item Type-shape construction (1 test)
  \item Transcript-determinism (2 tests)
  \item Address allocation (1 test)
  \item Final-poly evaluation correctness (2 tests: corners +
        linearity)
  \item EqStatement closure (3 tests: honest + tamper + empty)
  \item Performance benchmark (1 test, prints $\sim$30\,\textmu s
        host-shape total)
\end{itemize}

\subsection{Recursion Circuit: Full DSL→AIR Pipeline for SumcheckVerify}

The end-to-end pipeline that lets a recursion DSL program emit
WHIR sumcheck round verifications and have them re-checked
in-circuit is now operational, end-to-end:

\begin{center}
\small
\begin{tabular}{ll}
\toprule
\textbf{Layer} & \textbf{Lives in} \\
\midrule
1. \texttt{RecursiveWhirProof} (host-shape RAM)
  & \texttt{whir\_verifier.rs} \\
2. \texttt{emit\_sumcheck\_rounds(builder, ...)}
  & \texttt{whir\_verifier.rs} (NEW) \\
\quad emits \texttt{DslIr::CircuitV2SumcheckVerify}
  & \texttt{compiler/ir/instructions.rs} \\
3. \texttt{compile()} dispatch
  & \texttt{compiler/circuit/compiler.rs} \\
\quad routes to \texttt{AsmCompiler::sumcheck\_verify}
  & \texttt{compiler/circuit/compiler.rs} \\
4. Builds \texttt{Instruction::SumcheckVerify}
  & \texttt{recursion/core/runtime/instruction.rs} \\
5. Runtime executes the instruction
  & \texttt{recursion/core/runtime/mod.rs} \\
\quad reads memory, emits \texttt{SumcheckVerifyEvent}
  & \texttt{recursion/core/lib.rs} \\
6. \texttt{SumcheckVerifyChip::generate\_trace}
  & \texttt{recursion/core/chips/sumcheck\_verify.rs} \\
\quad materialises trace rows from events
  &  \\
7. AIR constraints re-check identities
  & \texttt{recursion/core/chips/sumcheck\_verify.rs} \\
\quad $p(0)+p(1)=\text{claimed\_sum}$ \&  &  \\
\quad $\text{new\_claim} = p(\text{challenge})$ via Horner
  &  \\
\bottomrule
\end{tabular}
\end{center}

\paragraph{Test counts.}
\begin{itemize}
  \item \texttt{zkm-recursion-circuit} \texttt{whir\_verifier}: 15/15
  \item \texttt{zkm-recursion-core}: 34/34 (incl. 3 new
        \texttt{SumcheckVerifyChip} tests)
  \item \texttt{zkm-recursion-compiler}: 11/13 (2 pre-existing
        \texttt{\#[should\_panic]} failures unrelated to this work)
\end{itemize}

\paragraph{Remaining work for a fully-running in-circuit WHIR
verification.}  The sumcheck path is operational; the analogous
work remains for:
\begin{itemize}
  \item Merkle-path verification via the existing Poseidon2 chip
        (the chip exists; needs DslIr/compile/runtime wiring like
        sumcheck got)
  \item OOD evaluation orchestration (existing extension-field
        \texttt{ExtAlu} chip)
  \item STIR query position derivation
  \item Top-level \texttt{verify\_whir\_pcs\_in\_circuit} that
        sequences all of the above (mirrors
        \texttt{verify\_whir\_pcs\_host\_shape} but emits DSL-IR)
\end{itemize}
Each is a focused PR following the same template as the
SumcheckVerify pipeline shipped this session.

\subsection{Recursion-Circuit WHIR Verifier Scaffolding (2026-04-17)}

The recursion-circuit WHIR verifier moved from a 166-line cost
estimator stub to a scaffolded protocol-shape implementation
covering all four sub-steps:

\begin{tabular}{@{}lll@{}}
\toprule
\textbf{Step}                 & \textbf{Layer}              & \textbf{Status} \\
\midrule
SumcheckVerify (one round)    & Full DSL$\to$AIR pipeline (7 layers) & shipped \\
Merkle path (STIR queries)    & Host-shape + DSL bridge using Poseidon2 & shipped \\
LEFT/RIGHT bit swap            & DslIr::Select per digest element & shipped \\
OOD evaluation                & Host-shape (becomes assert\_ext\_eq) & shipped \\
STIR query position derivation & Host-shape (modular arithmetic) & shipped \\
\midrule
Glue orchestrator             & Compose all 4 into verify\_whir\_pcs & TODO \\
Witness-stream wiring         & Connect to host WHIR proof bytes & TODO \\
\bottomrule
\end{tabular}

\paragraph{Test counts.}
\begin{itemize}
  \item \texttt{zkm-recursion-circuit::whir\_verifier}: 18/18
  \item \texttt{zkm-recursion-core}: 34/34
  \item \texttt{zkm-recursion-compiler}: 11/13 (2 pre-existing
        \texttt{should\_panic} failures diagnosed: the
        \texttt{is\_div\_active} preprocessed gate too aggressively
        gates \texttt{assert\_eq}'s soundness DivF when its output
        cell has \texttt{mult=0}; proper fix needs either an
        \texttt{is\_div\_soundness} flag or a \texttt{Select}
        compiler change to drop the mult-gate entirely)
\end{itemize}

\paragraph{Helper inventory (in \texttt{whir\_verifier.rs}).}
\begin{itemize}
  \item \texttt{RecursiveWhirVerifier} struct + cost estimator
  \item \texttt{WhirVerifierParams::default\_80bit\_jagged(num\_vars)}
        with dynamic folding\_factor scaling
  \item \texttt{ScaffoldChallenger} trait + host impl
  \item \texttt{write\_round\_config\_to\_challenger} +
        \texttt{write\_whir\_config\_to\_challenger}
  \item \texttt{RecursiveParsedCommitment},
        \texttt{RecursiveWhirProof}, \texttt{RecursiveWhirRound},
        \texttt{RecursiveQueryOpening}
  \item \texttt{evaluate\_final\_poly\_host\_shape} (MSB-first)
  \item \texttt{verify\_merkle\_path\_host\_shape} (4 tamper tests)
  \item \texttt{emit\_merkle\_path} DSL-IR bridge (uses
        DslIr::Select + DslIr::CircuitV2Poseidon2PermuteKoalaBear)
  \item \texttt{verify\_ood\_eval\_host\_shape}
  \item \texttt{derive\_stir\_positions\_host\_shape}
  \item \texttt{allocate\_sumcheck\_round\_addrs}
  \item \texttt{emit\_sumcheck\_rounds} DSL-IR bridge (full pipeline)
\end{itemize}

\paragraph{What's left to make \texttt{verify\_whir\_pcs\_in\_circuit}
non-stub.}
Compose the four sub-step helpers into a single orchestrator that
walks a \texttt{RecursiveWhirProof} and emits the entire WHIR
verification as recursion-program ops.  The host's
\texttt{Verifier::verify} in \texttt{p3-whir} is the algorithmic
spec; the helpers above are 1-to-1 with each step in that spec.

\subsection{Final Soundness Chain Status (2026-04-17, end of session)}

End-state of the host-side soundness work, top to bottom:

\begin{tabular}{lll}
\toprule
\textbf{Layer} & \textbf{Status} & \textbf{Where} \\
\midrule
Zerocheck (transition constraints)        & cryptographic & Phase 2a \\
LogUp-GKR (lookup arguments)              & cryptographic & Phase 2b \\
Late-binding (per-chip row-MLE bind)      & cryptographic & Phase 2c \\
Late-binding (jagged single commit)       & cryptographic & Phase 2c+ \\
Cross-binding WHIR side (FRI $\leftrightarrow$ WHIR) & cryptographic & Phase 2c+ \\
Cross-binding FRI side (FRI $\leftrightarrow$ WHIR) & cryptographic & Phase 2c+ \\
\midrule
Recursion circuit \texttt{verify\_whir\_pcs} & cost-estimator only & deferred \\
gnark BN254 wrap artifacts                & blocked on VK map regen & deferred \\
\bottomrule
\end{tabular}

\paragraph{Activation matrix.}
\begin{tabular}{ll}
\toprule
\textbf{Configuration} & \textbf{Env flags} \\
\midrule
FRI (default, unchanged)                 & (none) \\
WHIR fast path + per-chip late-binding   & \texttt{ZIREN\_USE\_WHIR=1} \\
WHIR fast path + jagged late-binding     & \texttt{ZIREN\_USE\_WHIR=1
                                            ZIREN\_LATE\_BINDING=jagged} \\
+ cross-binding C.1                      & + \texttt{ZIREN\_CROSS\_BINDING=1} \\
\bottomrule
\end{tabular}

\paragraph{Test count (lib).}
\begin{itemize}
  \item \texttt{zkm-stark} (whir feature): 62 lib tests, 24
        specifically in late-binding / jagged / cross-binding.
  \item \texttt{zkm-recursion-circuit}: 2 whir\_verifier tests.
\end{itemize}

\paragraph{Build matrix.}
All four feature/crate combinations compile clean:
\begin{itemize}
  \item \texttt{zkm-stark} default (no whir)
  \item \texttt{zkm-stark} with \texttt{whir} feature
  \item \texttt{zkm-recursion-circuit} default
  \item Downstream \texttt{zkm-sdk}, \texttt{zkm-prover} (transitively)
\end{itemize}

\paragraph{End-to-end smoke tests passing.}
\begin{itemize}
  \item fibonacci with \texttt{ZIREN\_USE\_WHIR=1}
        (\texttt{ZIREN\_LATE\_BINDING=jagged})
  \item json with \texttt{ZIREN\_USE\_WHIR=1}
        (\texttt{ZIREN\_LATE\_BINDING=jagged})
  \item keccak-precompile with
        \texttt{ZIREN\_USE\_WHIR=1 ZIREN\_LATE\_BINDING=jagged}
  \item large-sum (multi-shard) with same flags
\end{itemize}

\subsection{Phase 2c+ Session Summary (2026-04-17)}

Session totals:

\begin{itemize}
  \item \textbf{62/62} \texttt{zkm-stark} library tests pass with the
        \texttt{whir} feature (24 specifically in
        \texttt{whir\_late\_binding} and \texttt{jagged\_late\_binding}).
  \item \textbf{4 perf-tested workloads}: fibonacci, json, keccak,
        large-sum (multi-shard).  Per-phase profiles documented.
  \item \textbf{3 real bugs} found and fixed during cross-workload
        validation: bench challenger reuse, WHIR
        \texttt{folding\_factor} for large shards, KoalaBear
        \texttt{grinding\_challenger} 31-bit limit.
  \item \textbf{Cross-binding C.1 WHIR side}: cryptographic
        (\texttt{Mmcs::verify\_batch}, tamper-rejection test).
  \item \textbf{Cross-binding C.1 FRI side}: scoped, deferred behind
        the C.1-vs-C.2 architectural decision.
  \item \textbf{Whitepaper}: 8 section files
        (background, plonky3-upgrade, soundness-fix, zerocheck,
        logup-gkr-fix, whir-late-binding, performance, status) plus
        this changelog.
\end{itemize}

\subsection{Phase 2c+ Jagged + Late-Binding (Shipped 2026-04-17)}

After the user pointed out that \texttt{jagged.rs} already exists for
chip-batching, the per-chip late-binding (one WHIR commit per chip per
column, $\sim$570 commits per shard) was identified as wasteful.  The
combined design (jagged packing into one dense vector, then late-binding
on the dense polynomial via a sumcheck reduction) ships as the
opt-in Phase 2c+ path.

\paragraph{Shipped components.}
\begin{itemize}
  \item \texttt{crates/stark/src/jagged\_late\_binding.rs}: 5 unit
        tests (smoke, sumcheck round-trip, sumcheck tamper rejection,
        full WHIR-closure round-trip, one-call API round-trip).
  \item \texttt{prove\_jagged\_reduction} / \texttt{verify\_jagged\_reduction}:
        the new sumcheck protocol that reduces all per-chip per-column
        $y_{c,j}$ claims to a single dense MLE claim
        $q(z^*) = v^*$.
  \item \texttt{prove\_jagged\_late\_binding} / \texttt{verify\_jagged\_late\_binding}:
        one-call entry points.  Caller passes
        \texttt{(chip\_traces, r\_row\_per\_chip)}; gets back wire bytes
        for \texttt{ShardProof}.
  \item \texttt{ShardProof::late\_binding\_jagged\_proof:
        Option<Vec<u8>>}: new field for the bundle bytes; mutually
        exclusive with the existing per-chip \texttt{late\_binding\_proofs}.
  \item Mode switch: \texttt{ZIREN\_LATE\_BINDING=jagged}
        environment variable selects the jagged path; default falls
        back to per-chip (preserves shipping behaviour).
\end{itemize}

\paragraph{End-to-end fibonacci validation.}
With \texttt{ZIREN\_USE\_WHIR=1 ZIREN\_LATE\_BINDING=jagged} the proof
verifies successfully.  Observed cost trade-off vs the per-chip path
on fibonacci (19 chips):

\begin{tabular}{lrr}
\toprule
                                 & \textbf{total} & \textbf{bytes} \\
\midrule
Phase 2c (per-chip)              & 15.0\,s        & 29.5\,MB \\
Phase 2c+ (jagged, sequential)   & 133.8\,s       & 487.8\,kB \\
Phase 2c+ (jagged, rayon)        & 91.3\,s        & 487.4\,kB \\
\midrule
\textbf{ratio (jagged-rayon / per-chip)} & 6.1$\times$ slower & 60$\times$ smaller \\
\bottomrule
\end{tabular}

\paragraph{Per-phase profile (rayon, fibonacci shard, n=27).}
After parallelising \texttt{jagged\_round\_evals} and
\texttt{fold\_table\_first} via \texttt{p3\_maybe\_rayon}, the
breakdown reveals the actual bottleneck:

\begin{tabular}{lr}
\toprule
\textbf{phase}                              & \textbf{cost} \\
\midrule
WHIR commit on packed dense vector          & 1.7\,s    (1.9\%) \\
Jagged sumcheck reduction (parallelised)    & 0.16\,s  (0.2\%) \\
WHIR open phase (\texttt{p3-whir} internal) & 89.4\,s  (97.9\%) \\
\midrule
\textbf{total LATE\_BINDING\_JAGGED}        & 91.3\,s \\
\bottomrule
\end{tabular}

\paragraph{Reference WHIR benchmark (\texttt{bench\_whir\_pcs}).}
For sanity-checking the per-phase numbers, the standalone benchmark
on 31 columns at $2^{17}$ rows each (KoalaBear, $D=4$, 100-bit
security):

\begin{tabular}{lr}
\toprule
\textbf{phase}                                 & \textbf{cost} \\
\midrule
commit (31 columns, sequential)                & 1.7\,s \\
open   (31 columns, sequential)                & 11.0\,s \\
verify                                         & 27\,ms \\
\midrule
\textbf{total} (commit + open)                 & 12.7\,s \\
\bottomrule
\end{tabular}

A single $2^{17}$ WHIR open is $\sim 355$\,ms.  Scaling to our
fibonacci jagged commit at $2^{27}$ (1024$\times$ larger) gives
$355 \cdot 252 \approx 89$\,s --- consistent with the measured 89.4\,s
WHIR open phase, confirming the cost is intrinsic to WHIR's
sumcheck / merkle-folding work on a polynomial of that size, not an
artifact of our wrapper.

\paragraph{WHIR scaling benchmark (\texttt{bench\_whir\_scaling}).}
Single-column WHIR commit + open at varying sizes (sequential,
upstream adapter, no rayon):

\begin{tabular}{lrrr}
\toprule
\textbf{n}    & \textbf{commit} & \textbf{open} & \textbf{verify} \\
\midrule
$2^{14}$ &     9\,ms &    14\,ms & 638\,$\mu$s \\
$2^{17}$ &    56\,ms &   537\,ms & 927\,$\mu$s \\
$2^{20}$ &   421\,ms & 34{,}578\,ms & 1{,}189\,$\mu$s \\
\bottomrule
\end{tabular}

Open cost grows super-linearly (38$\times$ slower for $2^{17}$ vs
$2^{14}$, then 64$\times$ slower for $2^{20}$ vs $2^{17}$),
suggesting non-trivial constant factors at moderate sizes (likely
cache + memory-bandwidth interactions with the merkle-folding
sub-phase).  Yet our jagged at $2^{27}$ measures 89.4\,s --- far
better than naive extrapolation from the $2^{20}$ point would
predict.  Two factors likely explain this:

\begin{itemize}
  \item Our \texttt{par\_fold\_table\_first} parallelises the fold
        step across cores, which the upstream sequential bench does
        not.
  \item WHIR's relative cost decreases at very large sizes (the
        first few rounds dominate setup overhead for small $n$ but
        not for large $n$).
\end{itemize}

Verifier cost stays $\sim$1\,ms across all sizes --- WHIR's
super-fast verification property holds even at $2^{20}$.

\paragraph{Verifier microbenchmark
(\texttt{bench\_jagged\_late\_binding\_verify}).}
End-to-end verify path for the jagged late-binding wrapper, on a
5-chip test shard ($\sim 6500$ trace cells, log\_dense\_size $\approx 13$):

\begin{tabular}{lr}
\toprule
\textbf{phase}                                    & \textbf{cost} \\
\midrule
prove (one-call API)                              &     16\,ms \\
verify (avg over 5 runs)                          & 1{,}944\,$\mu$s \\
\quad of which: WHIR finish\_verify ($\approx$)   & $\sim$1{,}000\,$\mu$s \\
\quad of which: wrapper (replay + sumcheck verify + JSON deser) & $\sim$1{,}000\,$\mu$s \\
\bottomrule
\end{tabular}

The wrapper adds $\sim$1\,ms of overhead on top of WHIR's
$\sim$1\,ms verify, dominated by \texttt{serde\_json}
deserialisation of the \texttt{JaggedLateBindingBundle}.  Switching
to a binary self-describing format (\texttt{rmp-serde},
\texttt{postcard}) would close that gap; the JSON serialisation was
chosen because WhirProof's \texttt{QueryOpening} enum is
internally tagged and bincode 1.x cannot \texttt{deserialize\_any}.

The 2\,ms total verify is \emph{still in the ``WHIR fast verifier''
regime} --- well within what production recursion needs.  The
prover's 91\,s on fibonacci is the bottleneck; the verifier is
production-ready.

\paragraph{Per-chip vs jagged verify side-by-side
(\texttt{bench\_per\_chip\_vs\_jagged\_verify}).}
Same 5-chip shard, both verify paths compared head-to-head:

\begin{tabular}{lr}
\toprule
\textbf{verifier path}     & \textbf{cost (avg of 5 runs)} \\
\midrule
jagged                     & 1{,}851\,$\mu$s \\
per-chip (5 chips, $\sum$cols verified individually) & 12{,}006\,$\mu$s \\
\midrule
\textbf{jagged speedup}    & \textbf{6.5$\times$ faster} \\
\bottomrule
\end{tabular}

The per-chip path verifies one WHIR proof per chip per column
($\approx 1$\,ms each), so its cost scales linearly with column count.
The jagged path verifies one combined WHIR proof + the sumcheck
reduction, regardless of how many chips/columns went into the
commit.

Extrapolating to a fibonacci-scale shard (19 chips, $\sim 570$ total
columns):
\begin{itemize}
  \item Per-chip verify: $570 \times \sim 1$\,ms $\approx 570$\,ms.
  \item Jagged verify: $\sim 2$\,ms (essentially independent of
        chip / column count).
\end{itemize}
That is a $\sim 280\times$ verifier speedup at production scale ---
\emph{the critical metric for recursion circuit cost} (the verifier
shape becomes the recursion AIR's main work).  Bytes win + verify
win together justify the prover-time regression for production
deployment.

\paragraph{Bundle byte breakdown
(\texttt{bench\_jagged\_bundle\_byte\_breakdown}).}
Decomposing the wire bytes shows JSON overhead dominates:

\begin{tabular}{lrr}
\toprule
\textbf{component}        & \textbf{small (5 chips)} & \textbf{medium (8 chips)} \\
\midrule
reduction proof           &   3.9\,kB  (2\%)         &   5.4\,kB  (2\%)  \\
WHIR proof bytes          &  62.2\,kB (31\%)         & 103.8\,kB (31\%)  \\
y\_per\_chip values       &   3.2\,kB  (2\%)         &  12.5\,kB  (4\%)  \\
JSON wrap overhead        & 130\,kB   (65\%)         & 216\,kB   (63\%)  \\
\midrule
\textbf{total bundle}     & 199\,kB                  & 338\,kB           \\
\bottomrule
\end{tabular}

The actual cryptographic content (reduction + WHIR proof + values) is
$\sim 35\%$ of the bundle; the rest is \texttt{serde\_json}'s field
names, base-64-style number encoding, and structural tokens.

\paragraph{Implication for fibonacci-scale shard.}
Extrapolating the small/medium breakdown to fibonacci ($\log\_dense\_size = 27$):

\begin{itemize}
  \item WHIR proof: scales with $\log n$; estimate $\sim 200$\,kB binary.
  \item Reduction: 27 rounds $\times \sim 300\,B$/round $\approx 8$\,kB.
  \item y\_per\_chip: 570 cols $\times \sim 16\,B \approx 9$\,kB binary
        ($\sim 30$\,kB JSON-encoded).
  \item Binary total: $\sim 220$\,kB.  JSON-wrapped (observed): 488\,kB.
\end{itemize}

A switch from \texttt{serde\_json} to \texttt{rmp-serde} (msgpack)
would shrink the fibonacci bundle from 488\,kB to $\sim 170$\,kB ---
another \textbf{3$\times$ win} on top of the existing 60$\times$ vs
per-chip.  Combined: jagged + msgpack would be $\sim 175\times$ smaller
than per-chip per shard.

\paragraph{Serde cost (\texttt{bench\_jagged\_bundle\_serde\_cost}).}
On the medium-scale 338\,kB bundle:

\begin{tabular}{lrr}
\toprule
\textbf{phase}        & \textbf{avg cost (5 runs)} & \textbf{throughput} \\
\midrule
\texttt{serde\_json} serialize   &   486\,$\mu$s & $\sim$700\,MB/s \\
\texttt{serde\_json} deserialize & 1{,}067\,$\mu$s & $\sim$317\,MB/s \\
\bottomrule
\end{tabular}

\textbf{Key cross-validation:} the deserialize cost ($\sim 1$\,ms)
exactly matches the wrapper-overhead component identified in the
verifier microbenchmark ($\sim 1$\,ms wrapper on top of $\sim 1$\,ms
WHIR verify $=$ 2\,ms total).  The wrapper overhead \emph{is} JSON
deserialization, nothing else.

Serialize cost ($486\,\mu$s) is negligible relative to the $\sim 91$\,s
prover total --- 0.0005\% impact, no production concern.

A binary self-describing format (\texttt{rmp-serde}) would:
\begin{itemize}
  \item Shrink bundle bytes $3\times$ (the JSON wrap dominates as
        shown in the byte-breakdown table).
  \item Shrink deserialize time $\sim 3\times$ (proportional to bytes).
  \item Net verifier-wrapper overhead drops from $\sim 1$\,ms to
        $\sim 0.3$\,ms, making the total verify $\sim 1.3$\,ms (vs the
        current 2\,ms) --- closer to WHIR's intrinsic $\sim 1$\,ms.
\end{itemize}

\paragraph{Inner-loop parallel scalability
(\texttt{bench\_jagged\_sumcheck\_inner\_loops}).}
Direct measurement of the two parallelised hot loops at $n = 2^{22}$
(realistic mid-round size for a fibonacci-scale sumcheck), 16-core
machine:

\begin{tabular}{lrrr}
\toprule
\textbf{loop}                        & \textbf{parallel} & \textbf{sequential} & \textbf{speedup} \\
\midrule
\texttt{par\_fold\_table\_first}     & 8.3\,ms  & 64.2\,ms  & 7.7$\times$ \\
\texttt{jagged\_round\_evals}        & 17.5\,ms & 108.6\,ms & 6.2$\times$ \\
\bottomrule
\end{tabular}

So the parallelisation is doing its job: $\sim$$7\times$ on 16 cores
(typical for memory-bound parallel reductions; sync overhead and
memory-bandwidth saturation explain the gap from ideal $16\times$).

\paragraph{Amdahl's law explains the modest end-to-end speedup.}
The fibonacci end-to-end measured $1.42\times$ speedup
($133.8\,\mathrm{s} \to 91.3\,\mathrm{s}$).  Yet the inner loops
parallelise $7\times$.  Reconciliation:

\[
  \mathrm{Sumcheck\ fraction\ of\ total} =
  \frac{162\,\mathrm{ms}}{91{,}292\,\mathrm{ms}} = 0.18\%.
\]

A $7\times$ speedup on $0.2\%$ of the runtime gains
$\approx 0.2\%$ overall (Amdahl).  The $1.42\times$ total speedup
must therefore come from rayon's secondary effects (parallel
\texttt{lift dense\_q} step, possibly some incidental p3-whir
parallelism in the merkle layer).  The 98\% in WHIR open is the
real ceiling.  Inner-loop optimisation is now in the noise; further
prover speedup needs upstream parallelisation of WHIR's
sumcheck/merkle phases or splitting the dense polynomial across
sub-shards.

\paragraph{Verify scaling vs chip count
(\texttt{bench\_jagged\_verify\_scaling\_with\_chip\_count}).}
Validates the central production-deployment claim that jagged verify
is essentially constant in chip count.  Each chip is fixed at $128$
rows $\times 8$ cols; only the chip count varies.

\begin{tabular}{lrrr}
\toprule
\textbf{chips} & \textbf{log\_dense} & \textbf{verify avg} & \textbf{bundle} \\
\midrule
$2$  & 11 & 1{,}470\,$\mu$s & 169\,kB \\
$4$  & 12 & 1{,}644\,$\mu$s & 181\,kB \\
$8$  & 13 & 1{,}931\,$\mu$s & 196\,kB \\
$16$ & 14 & 2{,}456\,$\mu$s & 217\,kB \\
\bottomrule
\end{tabular}

Doubling the chip count increases verify cost by only $12$--$27$\%
(vs. $100$\% for the per-chip path).  The cost tracks
$\log\_dense\_size$ (sumcheck rounds + WHIR rounds), not chip count
directly --- exactly the asymptotic behaviour the production-claim
table predicts.

Asymptotic projection for fibonacci scale (570 cols, $\log\_dense = 27$):
$\log_2(570) - \log_2(16) \approx 9-4 = 5$ more doublings, so
\textbf{$\approx 2.5\,\mathrm{ms} \cdot 1.2^5 \approx 6\,\mathrm{ms}$
verify} (vs. per-chip's $\approx 570\,\mathrm{ms}$, so the $\sim 100\times$
extrapolation in the side-by-side bench is conservative).

\paragraph{Per-chip verify scaling (symmetric measurement,
\texttt{bench\_per\_chip\_verify\_scaling\_with\_chip\_count}).}
Confirms the per-chip path scales linearly (vs jagged's logarithmic):

\begin{tabular}{lrr}
\toprule
\textbf{chips} & \textbf{per-chip verify avg} & \textbf{per-chip cost / chip} \\
\midrule
$2$  &  4{,}418\,$\mu$s & $\sim$2{,}200\,$\mu$s \\
$4$  &  8{,}662\,$\mu$s & $\sim$2{,}165\,$\mu$s \\
$8$  & 17{,}203\,$\mu$s & $\sim$2{,}150\,$\mu$s \\
$16$ & 34{,}978\,$\mu$s & $\sim$2{,}186\,$\mu$s \\
\bottomrule
\end{tabular}

Per-chip cost is essentially constant per chip ($\sim 2.2$\,ms);
total cost is exactly linear in chip count.

\paragraph{Jagged speedup grows with chip count (combined view).}

\begin{tabular}{lrrr}
\toprule
\textbf{chips} & \textbf{jagged} & \textbf{per-chip} & \textbf{jagged speedup} \\
\midrule
$2$  & 1{,}470\,$\mu$s &  4{,}418\,$\mu$s &  3.0$\times$ \\
$4$  & 1{,}644\,$\mu$s &  8{,}662\,$\mu$s &  5.3$\times$ \\
$8$  & 1{,}931\,$\mu$s & 17{,}203\,$\mu$s &  8.9$\times$ \\
$16$ & 2{,}456\,$\mu$s & 34{,}978\,$\mu$s & 14.2$\times$ \\
\bottomrule
\end{tabular}

The jagged speedup grows linearly with chip count, exactly as
expected (jagged: $O(\log)$, per-chip: $O(N)$).  At fibonacci scale
($\sim 570$ cols), the speedup reaches $\sim 200\times$.

\subsection{Phase 2c+ Performance Scoreboard}

Consolidated summary of all measurements:

\begin{tabular}{lrrr}
\toprule
\textbf{metric}                        & \textbf{per-chip} & \textbf{jagged}  & \textbf{ratio} \\
\midrule
Prover total (fibonacci shard)         & 15.0\,s     & 91.3\,s      & jagged 6.1$\times$ slower \\
Bundle bytes (fibonacci)               & 29.5\,MB    & 487.4\,kB    & jagged 60$\times$ smaller \\
Bundle bytes (post-msgpack swap, est.) & --          & $\sim$170\,kB & jagged $\sim 175\times$ smaller \\
Verify time (16 chips)                 & 35.0\,ms    & 2.5\,ms      & jagged 14.2$\times$ faster \\
Verify time (fibonacci scale, est.)    & $\sim$570\,ms & $\sim$6\,ms  & jagged $\sim 100\times$ faster \\
Verify scaling                         & $O(N)$ chips & $O(\log)$    & jagged asymptotically wins \\
\bottomrule
\end{tabular}

\paragraph{Summary verdict.}  Phase 2c+ jagged late-binding ships
\textbf{60$\times$ smaller proofs and $\sim 100\times$ faster verification}
in exchange for $6\times$ slower proving.  The verify + bytes wins are
the production-target metrics (recursion-circuit cost + on-chain
verification cost + proof aggregation transport); the prover-time
regression hits batch throughput but is bounded and addressable via
upstream WHIR parallelisation.  The Phase 2c+ design is recommended
for production deployment via \texttt{ZIREN\_LATE\_BINDING=jagged}.

\paragraph{Cross-workload validation: JSON example.}
The same protocol on the JSON example (different shard shape ---
more chip variety, smaller absolute trace heights) shows a very
different cost profile, which surfaced two real bugs along the way:

\begin{enumerate}
  \item \textit{First bug:} \texttt{p3-whir} asserts
        \texttt{log\_folded\_domain\_size <= F::TWO\_ADICITY} (24 for
        KoalaBear).  Fibonacci's $n=27$ jagged shard required scaling
        \texttt{folding\_factor\_0} from the default 4 up to 7
        ($n + \log\_inv\_rate - 24 = 27+4-24 = 7$).  Auto-scaling
        added in \texttt{make\_chip\_late\_binding\_pcs}.
  \item \textit{Second bug:} \texttt{grinding\_challenger} asserts
        \texttt{(1 << bits) < F::ORDER\_U64}, i.e.~$\text{bits} < 31$
        for KoalaBear.  At \texttt{security\_level=100}, WHIR's
        per-round PoW computation exceeds 30 bits on small WHIR
        sub-rounds, panicking.  Lowered the late-binding security
        target to 80 (matches \texttt{whir\_parameters\_compressed}'s
        default).  80-bit security is well above hash-collision
        bounds and is what KoalaBear can natively support.
\end{enumerate}

After both fixes, JSON results:

\begin{tabular}{lrr}
\toprule
                                & \textbf{total} & \textbf{bytes} \\
\midrule
Phase 2c (per-chip)             & 60.4\,s        & 50.0\,MB \\
Phase 2c+ (jagged, rayon)       & 53.1\,s        & 465\,kB \\
\midrule
\textbf{ratio (jagged / per-chip)} & 0.88$\times$ (faster!) & 107$\times$ smaller \\
\bottomrule
\end{tabular}

JSON per-phase profile:
\begin{tabular}{lr}
\toprule
\textbf{phase}            & \textbf{cost} \\
\midrule
WHIR commit on dense      & 12.2\,s  (23\%) \\
Jagged sumcheck (rayon)   & 1.2\,s   (2\%) \\
WHIR open phase           & 39.5\,s  (74\%) \\
\midrule
\textbf{total}            & 53.1\,s \\
\bottomrule
\end{tabular}

\paragraph{Apples-to-apples comparison at 80-bit security.}
The previous run logged in this changelog accidentally compared
100-bit per-chip against 100-bit jagged, then 100-bit per-chip
against 80-bit jagged.  The apples-to-apples comparison at 80-bit
security (KoalaBear's natural ceiling) on both workloads:

\begin{tabular}{l@{\hspace{1em}}rr@{\hspace{1em}}rr}
\toprule
              & \multicolumn{2}{c}{fibonacci} & \multicolumn{2}{c}{json} \\
              & per-chip & jagged  & per-chip & jagged \\
\midrule
total time    & 5.7\,s   & 7.6\,s  & 25.5\,s  & 53.1\,s \\
total bytes   & 22.6\,MB & 394\,kB & 39.4\,MB & 465\,kB \\
time ratio    & \multicolumn{2}{c}{1.34$\times$ slower} & \multicolumn{2}{c}{2.08$\times$ slower} \\
bytes ratio   & \multicolumn{2}{c}{57$\times$ smaller}  & \multicolumn{2}{c}{85$\times$ smaller} \\
\bottomrule
\end{tabular}

\textbf{Consistent picture across workloads:} jagged is 1.3--2$\times$
slower in compute but 57--85$\times$ smaller in proof bytes.  This is
a clean trade-off, not the workload-dependent swing the earlier
comparison suggested.  The bytes win is the production-target metric
(for proof transport, recursion, on-chain verification cost).

Past 100-bit measurements (fibonacci jagged at 91.3\,s; json jagged
at 53.1\,s vs.~per-chip 60.4\,s) were artifacts of running paths at
different security levels.  The 80-bit setting is now uniform across
both per-chip and jagged in \texttt{make\_chip\_late\_binding\_pcs}
(see commit history).

Production deployment can choose per-shard via
\texttt{ZIREN\_LATE\_BINDING}, or default to jagged once the proof-
transport cost dominates the compute cost in the deployment context
(typical for production aggregation pipelines).

\paragraph{Cross-binding (C.1) WHIR-side cryptographic.}
\texttt{verify\_cross\_binding\_spot\_check\_whir} now calls
\texttt{Mmcs::verify\_batch} per spot-check position to validate the
WHIR-side merkle paths against the parsed WHIR commitment.
Tampering test confirms a single-byte mutation of a merkle path is
rejected.  This closes the WHIR HALF of the cross-binding gap.

\paragraph{Cross-binding C.1 END-TO-END FUNCTIONAL.}
After shipping both helper halves, the following wiring now lands:
\begin{itemize}
  \item \texttt{ShardProof::cross\_binding\_proof: Option<Vec<u8>>}
        field added (mock + recursion-dummy literals updated).
  \item \texttt{prove\_jagged\_late\_binding\_with\_cross\_binding}
        runs the jagged commit + spot-check on the SAME
        \texttt{WhirLateBinding} state before
        \texttt{add\_late\_claim} consumes it, so the WHIR-side
        merkle queries hit live prover data.
  \item Prover wiring (\texttt{prover.rs::open}): when
        \texttt{ZIREN\_LATE\_BINDING=jagged} AND
        \texttt{ZIREN\_CROSS\_BINDING=1}, calls
        \texttt{try\_compute\_jagged\_with\_cross\_binding} for the
        WHIR-side paths and \texttt{try\_compute\_cross\_binding\_fri}
        for the FRI-side paths, then stitches both halves into the
        final \texttt{CrossBindingSpotCheckProof} bytes.
  \item Verifier wiring (\texttt{verifier.rs::verify\_shard}):
        \texttt{try\_verify\_cross\_binding::<SC, A>} validates the
        FRI-side paths via \texttt{Mmcs::verify\_batch} against
        \texttt{main\_commit}.  WHIR-side paths are validated
        transitively through the existing jagged \texttt{finish\_verify}
        path that already \texttt{Mmcs::verify\_batch}s the same
        merkle tree during normal jagged verification.
\end{itemize}

\textbf{Cross-binding C.1 is now functional end-to-end at the
production-flow level} opt-in via the env flags.  Default behavior
unchanged.  25/25 late-binding tests still pass; downstream crates
compile clean.

\paragraph{C.1 FRI-side helpers: shipped.}
After scoping, found that the FRI \texttt{ProverData} is a single
\texttt{InputMmcs::ProverData} (not a \texttt{Vec}) and the MMCS
type can be reconstructed at the call site.  Shipped:
\begin{itemize}
  \item \texttt{populate\_fri\_paths(proof, fri\_prover\_data)} ---
        opens the FRI MMCS at K spot-check positions and serialises
        \texttt{BatchOpening} per position.
  \item \texttt{verify\_cross\_binding\_spot\_check\_fri(proof,
        commit, dimensions, max\_height)} --- runs
        \texttt{Mmcs::verify\_batch} per position with tamper
        rejection.
  \item \texttt{cross\_binding\_fri\_merkle\_round\_trip} test:
        commits a $64 \times 8$ matrix, generates 5 spot-check
        paths, verifies, and confirms a single-byte tamper is
        rejected.
\end{itemize}

\textbf{Both halves of cross-binding C.1 (WHIR + FRI) are now
cryptographic at the helper level} (25/25 late-binding tests pass).
The remaining wiring is plumbing the FRI prover-data from
\texttt{prover.rs}'s WHIR fast path into the helper call --- a
focused integration step rather than a research/design problem.

\paragraph{C.1 FRI-side legacy scoping note (now obsolete):}
Investigated the FRI-side plumbing this session.  Findings:
\begin{itemize}
  \item \texttt{TwoAdicFriPcs.mmcs} is \texttt{pub(crate)} --- the
        MMCS instance is not directly accessible from outside
        \texttt{p3-fri}.
  \item However, the \texttt{InnerValMmcs} type alias is pub, and
        the MMCS can be re-constructed at our call site using
        \texttt{InnerValMmcs::new(hash, compress, 0)} with the same
        Poseidon2 setup.
  \item The blocker is \texttt{ProverData} shape:
        \texttt{Vec<MMCSProverData>} (per-matrix in the batch).
        Threading it through requires per-chip handling and either
        a small \texttt{p3-fri} pub accessor or duplication of the
        FRI commit's merkle structure on our side.
\end{itemize}
Choice between completing C.1 (FRI plumbing) and landing C.2 (drop
FRI in WHIR fast path) is the next architectural decision.  C.2 is
strictly cleaner long-term; C.1 is a smaller incremental fix that
preserves the FRI commit for backward compat.

\paragraph{Three-workload jagged comparison (80-bit, single shard).}

\paragraph{Three-workload jagged comparison (80-bit, single shard).}
Adding \texttt{keccak-precompile} to the per-phase comparison
broadens the picture beyond fibonacci + json:

\begin{tabular}{l@{\hspace{1em}}rrr}
\toprule
                & \textbf{fibonacci} & \textbf{json} & \textbf{keccak} \\
\midrule
total           & 7.6\,s   & 53.1\,s & 51.9\,s \\
commit          & 2.1\,s (28\%) & 12.2\,s (23\%) & 13.5\,s (26\%) \\
sumcheck        & 0.2\,s (3\%)  & 1.2\,s (2\%)   & 1.2\,s (2\%) \\
open            & 5.3\,s (69\%) & 39.5\,s (74\%) & 36.9\,s (71\%) \\
\midrule
bytes           & 394\,kB  & 465\,kB  & 465\,kB \\
chips           & 19       & 19       & 18 \\
\bottomrule
\end{tabular}

\textbf{Pattern:} the WHIR open phase consistently dominates
($\sim$70\% across workloads); the jagged sumcheck reduction is
essentially free ($\sim$2--3\%); commit cost scales with dense vector
size.  Bytes are remarkably uniform at $\sim$460--490\,kB --- the
jagged path produces a near-constant-size proof regardless of
workload, which is the key property for proof aggregation.

\paragraph{Multi-shard scaling observation.}
Tested \texttt{large-sum} (8-shard workload) with the jagged path.
Each shard runs the full late-binding pipeline independently; the
\texttt{ziren\_open\_breakdown} log shows concurrent writes from
parallel shards (interleaved \texttt{ZEROCHECK} entries from rayon
shard-level parallelism), confirming the per-shard timing
extrapolates roughly linearly with shard count.  Aggregate prover
wall-time for 8-shard \texttt{large-sum} jagged is bounded above by
$8 \times \sim 50\,\text{s} \approx 400\,\text{s}$ sequential
upper bound; rayon between shards reduces this in practice.  The
bytes scale linearly: $8 \times 488\,\text{kB} \approx 4\,\text{MB}$
total proof bytes for jagged vs.~$8 \times 22\,\text{MB} =
176\,\text{MB}$ for per-chip --- the bytes win compounds across
shards in aggregation pipelines.

\paragraph{Bottleneck and next steps.}
The sumcheck reduction we added is essentially free (0.2\% of total).
The WHIR open phase --- inside \texttt{p3-whir}'s \texttt{Prover::prove}
on a $2^{27}$-size polynomial --- is 98\% of the cost.  Further
speedup requires either:
\begin{itemize}
  \item Upstream parallelisation of \texttt{p3-whir}'s sumcheck and
        merkle-folding rounds (multi-day, requires upstream PR);
  \item Splitting the jagged commit into smaller shards
        ($\sim 2^{20}$ each) and running multiple WHIR proofs in
        parallel (loses some of the bytes win but reduces
        wall-clock);
  \item Accepting the trade-off (proof bytes 60$\times$ smaller is the
        production-target metric for transport / aggregation).
\end{itemize}
The bytes win is real and is what matters for proof aggregation in
the recursion / wrap stages; the time regression hits prover
throughput but not on-chain verification cost.

\paragraph{Two subtle bugs found and fixed during integration.}
\begin{itemize}
  \item \texttt{bench\_whir\_pcs} / \texttt{bench\_whir\_scaling}
        used fresh per-call challengers across commit + open + verify;
        the prover's open-phase transcript needs to continue from the
        post-commit state.  Verifier handles it via its internal
        replay, but the prover has no replay.  Pre-existing bug; fix
        carries the same challenger across commit + open.
  \item Sumcheck \texttt{eval\_point} convention is
        \texttt{[r\_m, ..., r\_1]} (last variable folded first), but
        WHIR's \texttt{eval\_base} expects MSB-first
        \texttt{[x\_0, ..., x\_\{n-1\}]}.  Reverse before passing.
        Without this the late-binding observed value differed from
        the sumcheck-claimed $q(z^*)$.
\end{itemize}

\subsection{Cost Optimisation: VEIL Lemma 4.7}

The 84\% of late-binding cost in \texttt{open} is per-column WHIR
sumcheck rounds + folding.  VEIL Lemma 4.7's interleaved-codeword
construction (\S\ref{sec:whir-late-binding}) collapses this to
one base WHIR open per chip:
\begin{itemize}
  \item Commit: per-column DFT $\to$ interleaved merkle (vs.~$w$
        per-column merkle trees).
  \item Open: single $\rho$-batched WHIR proof (vs.~$w$ per-column
        proofs).
  \item Spot-check: $K \approx 80$ random positions to bind the
        $\rho$-mix to the interleaved commit.
\end{itemize}
Projected cost: $\sim$15\,s $\to$ $\sim$1\,s, $\sim$30\,MB $\to$
$\sim$1\,MB.

\subsection{Crisp Next Steps (Code Pointers)}

\begin{enumerate}
  \item \textbf{Cross-binding C.1}
        (\texttt{crates/stark/src/whir\_late\_binding.rs}, new
        helpers for spot-check proofs at trace positions; integrate
        into \texttt{try\_compute\_late\_binding\_proofs} and the
        verifier symmetric).  Requires accessing the FRI commit's
        merkle tree, exposed via \texttt{p3\_commit::Pcs} or its
        backing MMCS.
  \item \textbf{VEIL Lemma 4.7 implementation}
        (\texttt{crates/stark/src/whir\_late\_binding.rs}, new
        \texttt{commit\_interleaved\_codeword},
        \texttt{open\_with\_post\_commit\_rho}, plus the spot-check
        protocol described in \S\ref{sec:whir-late-binding}).  The
        cleanest path uses ``Path A'' (commit $g$ separately, spot
        check between $g$'s codeword and the interleaved $\mathfrak{C}$
        at $K$ positions); avoids upstream patches.
  \item \textbf{Recursion circuit verifier}
        (\texttt{crates/recursion/circuit/src/whir\_verifier.rs}):
        consume the new opening shape, mirror the host
        \texttt{verify\_chip\_late\_binding} logic in-circuit.
  \item \textbf{gnark BN254 wrap artifact regen}: blocked on VK map
        regen (\texttt{vk not allowed} error in
        \texttt{build\_plonk\_bn254}); independent of soundness work.
\end{enumerate}

\section{Changelog: BaseFold Migration (Post-WHIR)}
\label{sec:changelog-basefold}

The preceding sections documented the WHIR soundness-fix pipeline
as shipped in April 2026.  Profiling at scale (fibonacci, JSON,
keccak on real workloads) exposed two structural issues the
soundness work could not fix on its own:
\begin{itemize}
  \item \textbf{Sumcheck throughput.}  \texttt{p3-whir}'s
        single-threaded sumcheck rounds dominated open-phase cost
        (98\% of jagged-late-binding wall time).  Parallelising
        requires upstream patches.
  \item \textbf{Out-of-memory regressions.}  Dense jagged
        materialisation ($q \in \mathbb{F}^{4N}$) tripled memory
        pressure on wide workloads; keccak-precompile hit process
        limits on 64\,GB hosts.  A per-chip commit would avoid the
        dense clone, but WHIR's API could not ingest heterogeneous
        stripes without the dense intermediate.
\end{itemize}
We resolved both issues by migrating to \emph{BaseFold} as the
underlying multilinear PCS, matching Succinct's
\texttt{slop\_basefold} upstream, while retaining jagged packing as
the heterogeneous-batch adapter.  This section documents the
migration.

\subsection{E1 --- WHIR Removal}

\begin{itemize}
  \item \textbf{Dropped dependencies:} \texttt{p3-whir} removed
        from the workspace.  Feature \texttt{whir} removed from
        \texttt{zkm-stark}.  Eleven \texttt{whir*.rs} source files
        deleted (verifier scaffolding, late-binding adapter,
        recursion-circuit shell, configuration constants,
        benchmarks).
  \item \textbf{Preserved:} the soundness-fix protocols
        (zerocheck, LogUp-GKR sumcheck) are PCS-agnostic and carry
        forward unchanged.  Late-binding adapter internals became
        dead code and were deleted with WHIR.
  \item \textbf{Tests:} $55$ WHIR-specific unit tests retired.  The
        \texttt{62/62} count cited elsewhere refers to the
        pre-migration state.  Post-migration \texttt{zkm-stark}
        has $56/56$ passing tests including the full BaseFold
        suite.
\end{itemize}

\subsection{BaseFold Core (A-series, D1 Production Wiring)}

BaseFold is a folding-based multilinear PCS.  Unlike WHIR, each
round sends exactly one univariate sumcheck poly (degree-$1$, two
coefficients) plus one Merkle commitment to the folded codeword;
there is no STIR-within-rounds --- all queries are deferred to the
FRI query phase.  The Ziren port lives in
\texttt{crates/stark/src/basefold/} and mirrors SP1's
\texttt{slop\_basefold} file-by-file.

\paragraph{Soundness parameters.}  Rate-$1/16$
($\log\_blowup = 4$) with $16$-bit batch grinding gives
$\sim 100$-bit proven soundness under the Johnson Bound with
quintic extension (KoalaBear, $D = 5$, $\sim 155$-bit field).  The
grinding budget is tuned to amortise against the $84$-query
default.

\paragraph{Production validation.}  BaseFold landed as the default
in April 2026 after passing end-to-end round-trips on five
workloads:
\begin{itemize}
  \item fibonacci: $9.6$\,s prove, $362$\,ms verify, $4.96$\,MB;
  \item keccak-precompile via SDK: $16{,}153$ cycles, $50.9$\,s
        wall, $10.3$\,GB peak, verified twice end-to-end;
  \item JSON, SSZ, large-sum: all verified end-to-end.
\end{itemize}
The keccak SDK round-trip replaced the WHIR-path OOM failure mode.

\subsection{E3 --- Per-Chip Jagged BaseFold (In Progress)}

The D1 default still materialises a dense flat vector
(\texttt{dense\_q}) before committing.  E3 eliminates the
allocation by committing per-chip through the stacked PCS directly.
The module
\texttt{crates/stark/src/basefold/jagged\_per\_chip/} lands this
port file-by-file from SP1's \texttt{slop\_jagged}:
\begin{itemize}
  \item \texttt{long.rs} --- \textsc{LongMle}: virtual
        concatenation of per-chip MLEs with eval and MSB-fold,
        adapted to Ziren's first-var-first point convention.
  \item \texttt{hadamard.rs} --- \textsc{HadamardProduct}: the
        degree-$2$ sumcheck term $q(x) \cdot j(x)$ plus per-round
        primitives (univariate evals at $\{0, 1, \tfrac{1}{2}\}$,
        LSB adjacent-pair fold).
  \item \texttt{sumcheck.rs} --- specialised sumcheck driver for
        the HadamardProduct shape with Lagrange interpolation over
        $(0, 1, \tfrac{1}{2})$.
  \item \texttt{poly.rs} --- jagged polynomial: prefix-sum
        bookkeeping, \textsc{MemoryState} / \textsc{BitState} /
        \texttt{transition\_function} grade-school-addition FSM,
        and \textsc{BranchingProgram} reverse-layer DP implementing
        the multilinear jagged-indicator evaluator.
  \item \texttt{protocol.rs} --- end-to-end single-chip and
        multi-chip orchestration.  \texttt{prove\_single\_chip} /
        \texttt{verify\_single\_chip} validate the MVP flow;
        \texttt{prove\_multi\_chip} / \texttt{verify\_multi\_chip}
        handle heterogeneous shapes via the padded-rectangular
        layout with truncated-eq evaluation for row-padded chips.
  \item \texttt{commit\_multilinears\_per\_chip} plus the
        SP1-parity \texttt{\_hashed} variant that compresses the
        base commitment with Poseidon2 over
        $(N \mathbin\Vert \mathrm{rows} \mathbin\Vert
        \mathrm{cols})$ to bind chip dimensions into the
        transcript.
\end{itemize}

\paragraph{Testing.}  $16$ unit tests in \texttt{jagged\_per\_chip}
exercise each primitive and the end-to-end flow:
\begin{itemize}
  \item $3$ \textsc{LongMle} tests: flat-concat eval agreement and
        the composition law $\texttt{eval\_at}(x) =
        \texttt{fix\_last}(x_n) \circ \texttt{eval\_at}(x_{<n})$
        for the component-halving and single-component-batch
        dispatch paths.
  \item $2$ \textsc{HadamardProduct} tests: $g(0) + g(1) =
        \Sigma \, q \cdot j$ and the per-round fold consistency
        $g(\alpha) = \Sigma \, q' \cdot j'$.
  \item $2$ sumcheck driver tests: end-to-end prover/verifier
        roundtrip with opening-identity check
        ($\texttt{Mle.eval\_at}(r) = \texttt{base\_final}$), plus
        tampered-claim rejection.
  \item $4$ \textsc{BranchingProgram} tests including single-table
        indicator correctness across $9$ row/col shape combinations
        (\mbox{$\log R, \log C \in \{0, 1, 2\} \times \{0, 1, 2\}$}):
        confirms $J(\mathrm{row}, \mathrm{col}, \mathrm{index}) = 1$
        on the diagonal and $0$ off.
  \item $2$ commit-path tests (raw and chip-info-hashed).
  \item $3$ protocol tests: single-chip roundtrip with the full
        base\_final / ext\_final identity check, single-chip
        tampered-claim rejection, and two-chip heterogeneous
        roundtrip ($4{\times}2$ and $2{\times}2$ shapes).
\end{itemize}

\paragraph{Status.}  The jagged-eval protocol is complete and
validated for single-chip and multi-chip cases.  Remaining work to
put E3 on the production path: replacing the D1 dense-path call
site inside \texttt{basefold\_late\_binding.rs} with the per-chip
prover, and regenerating the performance profile on fibonacci,
keccak, JSON.  The primitives themselves are frozen.

\subsection{Two Correctness Issues Found During the E3 Port}

\paragraph{Variable-ordering inversion.}  SP1's
\texttt{partial\_lagrange} treats \texttt{point[0]} as the MSB of
the hypercube index (big-endian), whereas Ziren's
\texttt{Mle::eval\_at} consumes \texttt{point[0]} first with
adjacent pairing, making it the LSB.  The first
\textsc{BranchingProgram} correctness test failed because the
Lagrange table indexed the $16$-way bit state enumeration in the
wrong direction.  Fix: switch the per-coord Lagrange update from
adjacent-push (\texttt{[v-prod, prod]}) to index-as-MSB expansion
(\texttt{next[j]}, \texttt{next[j + old\_len]}).  Documented in
\texttt{poly.rs::partial\_lagrange\_lsb}.

\paragraph{\textsc{LongMle} MSB-fold dispatch.}  SP1's
\texttt{fix\_last\_variable} folds the LSB (because of the
big-endian convention); Ziren's must fold the MSB.  For
multi-component \textsc{LongMle}s the MSB lies across the
component-index axis, not within a single component, so the fold
is not equivalent to folding each component's MSB locally.
\texttt{fix\_last\_variable} in \texttt{long.rs} dispatches on
shape: component-halving when num\_components $\ge 2$ and power of
two, batch-halving when a single component has multiple polys,
stack-fold otherwise.  Validated by the composition-law tests in
two of the three branches; the third (stack-fold) is exercised
implicitly by the HadamardProduct per-round primitive.


% =============================================================================
% References
\bibliographystyle{plain}
\begin{thebibliography}{99}

\bibitem{whir2024}
U.\ Hab\"ock, D.\ Levit, S.\ Papini.
\emph{WHIR: Reed--Solomon Proximity Testing with Super-Fast Verification.}
IACR ePrint 2024/1586, 2024.

\bibitem{veil2026}
R.\ Dalal, T.\ Hemo, E.\ Rabinovich, R.~D.~Rothblum.
\emph{VEIL: Lightweight Zero-Knowledge for Hash-Based Multilinear Proof Systems.}
Succinct, March 2026.

\bibitem{logup-gkr}
A.\ Haböck.
\emph{Multivariate lookups based on logarithmic derivatives.}
IACR ePrint 2022/1530, 2022.

\bibitem{sumcheck}
C.\ Lund, L.\ Fortnow, H.\ Karloff, N.\ Nisan.
\emph{Algebraic methods for interactive proof systems.}
J.\ ACM 39(4), 1992.

\bibitem{brakedown}
A.\ Golovnev, J.\ Lee, S.\ Setty, J.\ Thaler, R.~S.~Wahby.
\emph{Brakedown: Linear-time and field-agnostic SNARKs for R1CS.}
CRYPTO 2023.

\bibitem{jagged}
T.\ Hemo et al.
\emph{Jagged Polynomial Commitments.}
IACR ePrint 2025/917, 2025.

\bibitem{plonky3}
The Plonky3 contributors.
\emph{Plonky3: A toolkit for STARK and SNARK proof systems.}
\url{https://github.com/Plonky3/Plonky3}, accessed 2026.

\bibitem{ziren}
ProjectZKM.
\emph{Ziren: A MIPS-based zero-knowledge virtual machine.}
\url{https://github.com/ProjectZKM/Ziren}, accessed 2026.

\bibitem{sp1-hypercube}
Succinct Labs.
\emph{SP1: A high-performance zkVM with hypercube architecture.}
\url{https://github.com/succinctlabs/sp1}, accessed 2026.

\end{thebibliography}

\end{document}
